\documentclass[../Thesis-IJspeert.tex]{subfiles}

\begin{document}

\graphicspath{ {"A Cavity-Compatible Setup for Single-Atom Traps/figs/"} }
\pgfplotsset{table/search path={"A Cavity-Compatible Setup for Single-Atom Traps/data/"}}

\chapter[A Cavity-Compatible Setup for Single-Atom Traps]{A Cavity-Compatible Setup for\\Single-Atom Traps}
\addtocontents{toc}{\vskip-6pt\par\noindent\protect\textcolor{gray75}{\protect\rule{\textwidth}{0.5pt}}\par}
\label{chap:CavityCompatibleTrappingArchitectures}
In this chapter, we present the design and implementation of a mechanical assembly comprising a high-finesse optical resonator, a tightly focused optical dipole trap and a compact cold-atom source based on a pyramidal magneto-optical trap (MOT). Following an introduction that sets the broader context for cavity-compatible atom trapping, we discuss the general mounting structure in \autoref{sec:mounting_structure}, where we first provide an overview of its geometry and integration constraints, before presenting a detailed mechanical stability and vibration analysis to ensure reliable operation in an ultra-high vacuum environment. In \autoref{sec:pyramidal_mot}, we describe the pyramidal MOT, starting from its basic operating principle and continuing with an analysis of aperture diffraction effects, the fabrication of the right-angle prism mirrors and the methods used for cold-atom detection. The chapter closes with a summary of the main design choices and implications for future work on intra-cavity single-atom trapping.
\section{Introduction}
Efficient atom-photon interfaces are a key enabling technology for quantum information processing. In particular, their potential for deterministic single-photon generation and the distribution of remote entanglement across a quantum network. As discussed in \autoref{chap:TheoryandDesignofanOpticalResonator}, the atom-photon coupling can be strengthened significantly by placing the emitter in a high-finesse optical cavity. This supports deterministic photon emission into a single mode of the radiation field through what is known as the Purcell effect \cite{Purcell1946}. The potential of a cavity-enhanced atom-photon interface for the delivery of single photons and distributed entanglement has been demonstrated using an atomic fountain, by means of which \ce{^87Rb} atoms are stochastically loaded into a high-finesse optical cavity \cite{Barrett2019}. In this case, the loading efficiency is deliberately kept low to suppress the likelihood of multi-atom occupancy within the resonator, which limits the interaction time and hence the photon generation efficiency. This is further limited by the transient nature of atomic loading in this experiment, which gives rise to time-dependent coupling strengths and a finite time frame for the repeated execution of V-STIRAP. In addition, stochastic loading becomes ineffective when multiple cavities must be loaded simultaneously---a requirement for realising entanglement protocols based on the inference of indistinguishable photons emitted from separate atom-cavity systems (see \autoref{app:cavityentanglement}). This follows from the fact that the likelihood of loading two cavities at the same time is comparable to the likelihood of finding two atoms within the mode volume of a single resonator (see \autoref{app:stochasticcavityloading}).

While such platforms constitute excellent test beds, probabilistic loading fundamentally limits scalability. To maximise the single-atom occupation probability of an optical resonator, an atom can be confined in a tightly focused optical dipole trap operating in the collisional blockade regime (see \autoref{collisionalblockade}). By employing holographic optical tweezer arrays, multiple atoms can be positioned with high precision within the cavity mode, enabling simultaneous coupling of several emitters. Alternatively, such arrays can act as a quantum register, where individual atoms are stored outside the mode volume and dynamically shuttled in and out to interact with the cavity field as required.

One challenge with this approach, however, is that direct loading of the dipole traps from a MOT is often impractical due to the restricted optical access for the MOT beams in the space between the cavity mirrors. Instead, the MOT cloud must be positioned in the vicinity of the cavity, from which atoms are subsequently transported into the trapping region inside the resonator. A number of such approaches have been implemented, including funnelling atoms from an atomic fountain into the cavity using blue-detuned cavity modes \cite{Puppe2007}, guiding atoms from a static MOT with a single red-detuned beam \cite{Numann2005} and employing an optical lattice as a conveyor belt for atom transport \cite{Gallego2018,Schrader2001}.

Here, our approach for a cavity-compatible cold-atom source is based on the pyramidal MOT concept first introduced in 1996 \cite{Lee1996}. In this geometry, four reflective facets form a hollow pyramid that redirects a single incident cooling beam to create a compact three-dimensional trap. When configured with a small aperture at the apex, this setup can produce a continuous beam of cold atoms emanating along the pyramid axis \cite{Arlt1998,Williamson1998}. When directed towards the trapping region between the cavity mirrors, the atomic beam serves to load the intra-cavity dipole traps. In the sections that follow, we examine the design choices and engineering considerations underpinning a modular apparatus that integrates the optical resonator presented in \autoref{chap:TheoryandDesignofanOpticalResonator}, a pyramidal MOT cold-atom source and single-atom dipole traps.

\section{Mounting structure}
\label{sec:mounting_structure}
The central component of the cavity-compatible trapping setup is a rigid mounting structure which secures the cavity module whilst also accommodating the in-vacuum optics for the pyramidal MOT source inside a glass cell. This section outlines the design principles and overall layout of the assembly and evaluates its stability through detailed mechanical and vibrational analysis, with reference to an earlier mounting design for a fibre-tip optical resonator.

\subsection{Mechanical stability of a Fabry--Pérot resonator}
\label{subsec:mechanicalstab}
Fabry--Pérot resonators underpin a wide range of applications, from cavity optomechanical sensing \cite{Aspelmeyer2014} and gravitational-wave detection \cite{Abbott2016} to optical clocks \cite{Hinkley2013} and CQED experiments \cite{Mabuchi2002,Birnbaum2005,Steiner2013}. In all these contexts, active stabilisation of the cavity length is essential to suppress environmental fluctuations and maintain resonance. This is achieved by locking to an external frequency reference via feedback to the actuators that control the mirror positions. The usable feedback bandwidth is ultimately constrained by the mechanical response of the resonator assembly, since sharp vibrational resonances in the mirrors or their mounts directly affect the control loop and reduce stability.

A variety of approaches have been explored to improve the mechanical stability of Fabry--Pérot cavities, especially in the fields of laser frequency stabilisation and interferometry \cite{Notcutt2005,Briles2010,Sinclair2015}. For CQED experiments, however, the constraints of compact UHV environments limit the direct use of such damping or mounting strategies. Active suppression methods, for instance employing digital filters to cancel mechanical resonances within feedback loops \cite{Ryou2017}, can provide improvements, but add considerable complexity to the stabilisation system. A more robust solution is to design the cavity assembly for passive stability from the outset, reducing the presence of low-frequency resonances that would otherwise need to be actively compensated. Previous work has shown that resonances can be shifted to higher frequencies and suppressed in quality factor by maximising stiffness and introducing damping in the mounting structure \cite{Briles2010, Sinclair2015}. For the glass-substrate Fabry--Pérot resonator introduced in \autoref{chap:TheoryandDesignofanOpticalResonator}, we pursue passive stability by maximising the steel mass that can be accommodated within the compact glass cell and by introducing geometric irregularities in the mount to avoid degenerate mechanical modes.

Before presenting the design and modal frequency response analysis of this mount, we first examine an earlier-generation fibre-tip resonator and its corresponding mount. The following three sections characterise this fibre-tip system through simulated modal frequency response, assessment of passive stability under environmental noise and measurement of the actuator transfer function, providing a basis for comparison.

\subsection{Modal frequency response analysis of a fibre-tip resonator mount}
\label{modalfreqresponseanalysisfibretip}
Modal frequency response analysis is a computational technique used to predict how a structure responds dynamically to base excitations across a range of frequencies (see \autoref{app:modalanalysis}). By decomposing the system into its eigenmodes, each characterised by a natural frequency and mode shape, the method allows the evaluation of the contribution of individual modes to the overall dynamic behaviour. This approach is particularly effective for multi-degree-of-freedom (MDOF) systems, where the number of vibrational modes can be large, but only a subset significantly influences the response under typical operating conditions.

The modal frequency response of the fibre-tip resonator mount is evaluated through finite-element modelling, implemented in Autodesk Inventor using the Nastran solver. The analysis uses \num{8} processor threads\footnote{Using around \num{8} threads is near optimal for this analysis. Using more threads increases overhead from distributing the calculation and does not necessarily reduce total runtime.} on an AMD Ryzen Threadripper 3990X 64-Core Processor (\SI{2.90}{\giga\hertz}) with \SI{256}{\giga\byte} of RAM. A parabolic mesh with an element size of \SI{1}{\milli\metre} is used, resulting in a total of \num{85527} elements. To capture the dominant dynamic behaviour, \num{80} modes are retained in the analysis, such that the cumulative effective modal mass reaches at least \SI{80}{\percent} of the total physical mass for each degree of freedom. The effective modal mass ratios are listed in \autoref{modeparticipationM} (see \autoref{app:effectivemodalmass}). Base excitations with an acceleration magnitude of \SI{1}{\milli\metre\per\second\squared} are enforced individually in the $x$, $y$ and $z$ directions at the mounting interface to the vacuum chamber flange. The frequency response is evaluated between \SI{0}{\kilo\hertz} and \SI{10}{\kilo\hertz} with \num{15} data points between each pair of successive modes and a bias factor of \num{3} to ensure denser sampling near the resonance frequencies. The maximum linear acceleration magnitudes of the structure for these three excitation directions are presented in \autoref{freqresponseMarw}.
\begin{figure}[!t]
	%\vspace{0.7em}
	\centering
	\begin{tikzpicture}
	
	\begin{groupplot}[group style={group size=1 by 3, vertical sep=0.5cm},xmin=0,ymin=0,height=4cm,width=13.5cm,no markers]
	\nextgroupplot
	[
	ymode=log,
	ymin=1e-0,
	ymax=1e4,
	%xlabel={Frequency},
	%ylabel={Max. linear acceleration},
	xmin=0, xmax=10,
	xticklabels={,,},
	%ymin=-0.1, ymax=1.1,
	%ymin=0, ymax=120,
	%xtick={0,20,40,60,80,100},
	%ytick={0,20,40,60,80,100,120},
	%legend pos=north west,
	ymajorgrids=true,
	xmajorgrids=true,
	grid style=dashed,
	legend cell align={left},
	legend style={
		at={(1, 1)},  % Adjust (x, y) to move it
		anchor=north east,  % Attach the legend to its top-right corner
		draw=none,           % Optional: Remove border around the legend
		fill=none,
		align=left
	}
	]
	
	%\node[align=center] at (axis cs: 0.267, 3500) {\footnotesize A};
	\node[align=center] at (axis cs: 1.105, 3500) {\footnotesize\sc b};
	\node[align=center] at (axis cs: 1.254+0.2, 11) {\footnotesize\sc c};
	\node[align=center] at (axis cs: 2.940, 2000) {\footnotesize\sc d};
	%\node[align=center] at (axis cs: 3.024, 2500) {\footnotesize E};
	\node[align=center] at (axis cs: 3.893, 1500) {\footnotesize\sc f};
	%\node[align=center] at (axis cs: 6.643, 300) {\footnotesize G};
	\node[align=center] at (axis cs: 6.935, 1600) {\footnotesize\sc h};
	
	
	\addplot[
	color=markrood,
	name path=xpointone,
	] table [col sep = comma, x expr = \thisrowno{0}/1000, y index = 1]{M_x_01.csv};
	
	\addplot[
	color=rainbow2of8,
	name path=xone,
	] table [col sep = comma, x expr = \thisrowno{0}/1000, y index = 1]{M_x_1.csv};
	
	\addplot[
	color=rainbow4of8,
	name path=xfive,
	] table [col sep = comma, x expr = \thisrowno{0}/1000, y index = 1]{M_x_5.csv};
	
	\addplot+[draw=none,domain=0:10,name path=xzero] {0.01}; 
	
	\addplot[color=markrood, opacity=0.5] fill between[
	of = xone and xpointone,
	]; 
	
	\addplot[color=rainbow2of8, opacity=0.5] fill between[
	of = xfive and xone,
	]; 
	
	\addplot[color=rainbow4of8, opacity=0.5] fill between[
	of = xzero and xfive,
	]; 
	
	\node at (rel axis cs: 0.52, 0.6) [anchor=center] {\includegraphics[width=3.5cm]{MARWANFINGERtrans.png}};
	
	\addlegendentry{\footnotesize \hspace{0.3em} $G=\SI{0.1}{\percent}$}
	\addlegendentry{\footnotesize \hspace{0.3em} $G=\SI{1}{\percent}$}
	\addlegendentry{\footnotesize \hspace{0.3em} $G=\SI{5}{\percent}$}
	
	\draw[<->] (rel axis cs: 0.64, 0.75) --++ (150:0.5cm);
	\node[align=left] at (rel axis cs: 0.632, 0.88) {\footnotesize $x$};
	
	\nextgroupplot
	[
	ymode=log,
	ymin=1e-0,
	ymax=1e3,
	%ymin=1e-2, ymax=1e0,
	%xlabel={Frequency},
	ylabel={Maximum linear acceleration magnitude {[\si{\milli\meter\per\second\squared}]}},
	xmin=0, xmax=10,
	xticklabels={,,},
	%ymin=-0.1, ymax=1.1,
	%ymin=0, ymax=120,
	%xtick={0,20,40,60,80,100},
	%ytick={0,20,40,60,80,100,120},
	%legend pos=north west,
	ymajorgrids=true,
	xmajorgrids=true,
	grid style=dashed,
	legend cell align={left},
	legend style={
		at={(1, 1)},  % Adjust (x, y) to move it
		anchor=north east,  % Attach the legend to its top-right corner
		draw=none,           % Optional: Remove border around the legend
		fill=none,
		align=left
	}
	]
	
	\node[align=center] at (axis cs: 0.267, 2.8) {\footnotesize\sc a};
	%\node[align=center] at (axis cs: 1.105, 3500) {\footnotesize B};
	\node[align=center] at (axis cs: 1.254, 33) {\footnotesize\sc c};
	\node[align=center] at (axis cs: 2.940-0.2, 2) {\footnotesize\sc d};
	\node[align=center] at (axis cs: 3.024, 43) {\footnotesize\sc e};
	\node[align=center] at (axis cs: 3.893, 2.5) {\footnotesize\sc f};
	\node[align=center] at (axis cs: 6.643, 320) {\footnotesize\sc g};
	%\node[align=center] at (axis cs: 6.935, 300) {\footnotesize H};
	
	\addplot[
	color=markrood,
	name path=ypointone,
	] table [col sep = comma, x expr = \thisrowno{0}/1000, y index = 1]{M_y_01.csv};
	
	\addplot[
	color=rainbow2of8,
	name path=yone,
	] table [col sep = comma, x expr = \thisrowno{0}/1000, y index = 1]{M_y_1.csv};
	
	\addplot[
	color=rainbow4of8,
	name path=yfive,
	] table [col sep = comma, x expr = \thisrowno{0}/1000, y index = 1]{M_y_5.csv};
	
	\addplot+[draw=none,domain=0:10,name path=yzero] {0.01}; 
	
	\addplot[color=markrood, opacity=0.5] fill between[
	of = yone and ypointone,
	]; 
	
	\addplot[color=rainbow2of8, opacity=0.5] fill between[
	of = yfive and yone,
	]; 
	
	\addplot[color=rainbow4of8, opacity=0.5] fill between[
	of = yzero and yfive,
	]; 
	
	\node at (rel axis cs: 0.52, 0.6) [anchor=center] {\includegraphics[width=3.5cm]{MARWANFINGERtrans.png}};
	
	%\addlegendentry{\small \hspace{0.3em} $G=\SI{0.1}{\percent}$}
	%\addlegendentry{\small \hspace{0.3em} $G=\SI{1}{\percent}$}
	%\addlegendentry{\small \hspace{0.3em} $G=\SI{5}{\percent}$}
	
	\draw[<->] (rel axis cs: 0.605, 0.76) --++ (30:0.5cm);
	\node[align=left] at (rel axis cs: 0.615, 0.895) {\footnotesize $y$};
	
	\nextgroupplot
	[
	ymode=log,
	ymin=1e-0,
	ymax=1e4,
	%ymin=1e-2, ymax=1e0,
	xlabel={Enforced acceleration frequency {[\si{\kilo\hertz}]}},
	%ylabel={Max. linear acceleration},
	xmin=0, xmax=10,
	%ymin=-0.1, ymax=1.1,
	%ymin=0, ymax=120,
	%xtick={0,20,40,60,80,100},
	%ytick={0,20,40,60,80,100,120},
	%legend pos=north west,
	ymajorgrids=true,
	xmajorgrids=true,
	grid style=dashed,
	legend cell align={left},
	legend style={
		at={(1, 1)},  % Adjust (x, y) to move it
		anchor=north east,  % Attach the legend to its top-right corner
		draw=none,           % Optional: Remove border around the legend
		fill=none,
		align=left
	}
	]
	
	\node[align=center] at (axis cs: 0.267, 3500) {\footnotesize\sc a};
	%\node[align=center] at (axis cs: 1.105, 3500) {\footnotesize B};
	\node[align=center] at (axis cs: 1.254, 2200) {\footnotesize\sc c};
	\node[align=center] at (axis cs: 2.940-0.2, 30) {\footnotesize\sc d};
	\node[align=center] at (axis cs: 3.024, 1600) {\footnotesize\sc e};
	\node[align=center] at (axis cs: 3.893, 4.5) {\footnotesize\sc f};
	\node[align=center] at (axis cs: 6.643, 1500) {\footnotesize\sc g};
	%\node[align=center] at (axis cs: 6.935, 300) {\footnotesize H};
	
	\addplot[
	color=markrood,
	name path=zpointone,
	] table [col sep = comma, x expr = \thisrowno{0}/1000, y index = 1]{M_z_01.csv};
	
	\addplot[
	color=rainbow2of8,
	name path=zone,
	] table [col sep = comma, x expr = \thisrowno{0}/1000, y index = 1]{M_z_1.csv};
	
	\addplot[
	color=rainbow4of8,
	name path=zfive,
	] table [col sep = comma, x expr = \thisrowno{0}/1000, y index = 1]{M_z_5.csv};
	
	\addplot+[draw=none,domain=0:10,name path=zzero] {0.01}; 
	
	\addplot[color=markrood, opacity=0.5] fill between[
	of = zone and zpointone,
	]; 
	
	\addplot[color=rainbow2of8, opacity=0.5] fill between[
	of = zfive and zone,
	]; 
	
	\addplot[color=rainbow4of8, opacity=0.5] fill between[
	of = zzero and zfive,
	]; 
	
	\node at (rel axis cs: 0.52, 0.6) [anchor=center] {\includegraphics[width=3.5cm]{MARWANFINGERtrans.png}};
	
	\draw[<->] (rel axis cs: 0.615, 0.74) --++ (90:0.5cm);
	\node[align=left] at (rel axis cs: 0.635, 0.86) {\footnotesize $z$};
	
	
	\end{groupplot}
	
	
	
	\end{tikzpicture}
	\caption[Modal frequency response analysis of the fibre-tip resonator mounting structure]{Modal frequency response analysis of the fibre-tip resonator mounting structure. The graphs show the maximum linear acceleration magnitude of the structure versus the frequency of enforced acceleration at the base in three different directions: $x$ \textsc{(top)}, $y$ \textsc{(middle)} and $z$ \textsc{(bottom)}. In each direction, traces are shown for three different values of the structural damping coefficient $G$. The magnitude of the enforced acceleration equals \SI{1}{\milli\meter\per\second\squared}}
	\label{freqresponseMarw} 
\end{figure}
Three values (\SI{0.1}{\percent}, \SI{1}{\percent} and \SI{5}{\percent}) of the structural damping coefficient $G$ are considered, representing different levels of material damping. For a modal frequency response analysis in Autodesk Nastran, structural damping is defined in units of $G = 2\zeta$, where $\zeta$ is the critical damping ratio. Consequently, the effective modal damping applied is $\sfrac{G}{2}$. Literature values for steel typically report a critical damping ratio below \SI{1}{\percent} \cite[p.~500]{Adams1998}, which corresponds to $G < \SI{2}{\percent}$ in this convention.
\startsquarepar
The modal frequency response analysis in \autoref{freqresponseMarw} shows that the mounting structure exhibits eight vibrational modes below \SI{10}{\kilo\hertz}. The corresponding mode shapes, shown as exaggerated deformations at representative amplitude, are illustrated in \autoref{marwanmodes}. The fundamental mode at \SI{267}{\hertz} \textsc{(a)} corresponds to a first-order cantilever bending mode of the mounting arm. This mode is predominantly excited by vertical base acceleration (along the $z$-axis), while its amplitude under excitation along the $x$- and $y$-axes is comparatively minor. A distinction can be made between the even modes \textsc{(b, d, f, h)} and the odd modes \textsc{(a, c, e, g)}. The even modes are primarily excited by base acceleration along the $x$-axis, whereas the odd modes are predominantly excited along the $z$-axis and, to a lesser extent, the $y$-axis. This distinction is practically relevant because the resonator fibres are mounted on either side of the access hole for the dipole-trap beam, where structural
\stopsquarepar
\begin{figure}[H]
	\vspace{0em}
	\centering
	\begin{subfigure}[b]{.45\linewidth}
		\begin{tikzpicture}
		\centering
		\node[inner sep=0] (image) at (0,0) {\includegraphics[width=5cm] {1F.png}};;
		\begin{scope}[x={(image.south east)},y={(image.north west)}]
		\node at (0.3, 0.7) {\footnotesize \SI{267}{\hertz}};
		\end{scope}
		\end{tikzpicture}
		\centering
		\caption{}
	\end{subfigure}
	\begin{subfigure}[b]{.45\linewidth}
		\begin{tikzpicture}
		\centering
		\node[inner sep=0] (image) at (0,0) {\includegraphics[width=5cm] {2F.png}};;
		\begin{scope}[x={(image.south east)},y={(image.north west)}]
		\node at (0.3, 0.7) {\footnotesize \SI{1105}{\hertz}};
		\end{scope}
		\end{tikzpicture}
		\centering
		\caption{}
	\end{subfigure}
	\\
	\begin{subfigure}[b]{.45\linewidth}
		\begin{tikzpicture}
		\centering
		\node[inner sep=0] (image) at (0,0) {\includegraphics[width=5cm] {3F.png}};;
		\begin{scope}[x={(image.south east)},y={(image.north west)}]
		\node at (0.3, 0.7) {\footnotesize \SI{1254}{\hertz}};
		\end{scope}
		\end{tikzpicture}
		\centering
		\caption{}
	\end{subfigure}
	\begin{subfigure}[b]{.45\linewidth}
		\begin{tikzpicture}
		\centering
		\node[inner sep=0] (image) at (0,0) {\includegraphics[width=5cm] {4F.png}};;
		\begin{scope}[x={(image.south east)},y={(image.north west)}]
		\node at (0.3, 0.7) {\footnotesize \SI{2940}{\hertz}};
		\end{scope}
		\end{tikzpicture}
		\centering
		\caption{}
	\end{subfigure}
	\\
	\begin{subfigure}[b]{.45\linewidth}
		\begin{tikzpicture}
		\centering
		\node[inner sep=0] (image) at (0,0) {\includegraphics[width=5cm] {5F.png}};;
		\begin{scope}[x={(image.south east)},y={(image.north west)}]
		\node at (0.3, 0.7) {\footnotesize \SI{3024}{\hertz}};
		\end{scope}
		\end{tikzpicture}
		\centering
		\caption{}
	\end{subfigure}
	\begin{subfigure}[b]{.45\linewidth}
		\begin{tikzpicture}
		\centering
		\node[inner sep=0] (image) at (0,0) {\includegraphics[width=5cm] {6F.png}};;
		\begin{scope}[x={(image.south east)},y={(image.north west)}]
		\node at (0.3, 0.7) {\footnotesize \SI{3893}{\hertz}};
		\end{scope}
		\end{tikzpicture}
		\centering
		\caption{}
	\end{subfigure}
	\\
	\begin{subfigure}[b]{.45\linewidth}
		\begin{tikzpicture}
		\centering
		\node[inner sep=0] (image) at (0,0) {\includegraphics[width=5cm] {7F.png}};;
		\begin{scope}[x={(image.south east)},y={(image.north west)}]
		\node at (0.3, 0.7) {\footnotesize \SI{6643}{\hertz}};
		\end{scope}
		\end{tikzpicture}
		\centering
		\caption{}
	\end{subfigure}
	\begin{subfigure}[b]{.45\linewidth}
		\begin{tikzpicture}
		\centering
		\node[inner sep=0] (image) at (0,0) {\includegraphics[width=5cm] {8F.png}};;
		\begin{scope}[x={(image.south east)},y={(image.north west)}]
		\node at (0.3, 0.7) {\footnotesize \SI{6935}{\hertz}};
		\draw[->] (1, 0.15) --++ (90:0.5cm) node[above, anchor=south] {\footnotesize $z$};
		\draw[->] (1, 0.15) --++ (150:0.5cm) node[above, anchor=base] {\footnotesize  $x$\,\, \phantom{x}};
		\draw[->] (1, 0.15) --++ (210:0.5cm) node[above, anchor=mid] {\footnotesize $y$\,\, \phantom{x}};
		\end{scope}
		\end{tikzpicture}
		\centering
		\caption{}
	\end{subfigure}
	\caption[Normal modes below \SI{10}{\kilo\hertz} of the fibre-tip resonator mounting structure]{Normal modes and corresponding frequencies below \SI{10}{\kilo\hertz} of the fibre-tip resonator mounting structure. The colour scale corresponds to the displacement from the unperturbed position of the structure (dashed contours). Mode labels correspond to those in \autoref{freqresponseMarw}.}
	\label{marwanmodes}
\end{figure}
deformation directly alters the cavity length. The fundamental mode already affects the cavity length through this effect, but higher-order odd modes introduce stronger local curvature of the mounting arm at the fibre position, leading to much larger cavity-length fluctuations.

Further insight is gained by comparing the effective modal mass across each translational and rotational degree of freedom (\autoref{modeparticipationM}). The cumulative effective modal mass is lowest for both translational ($T_y$) and rotational motion ($R_y$) associated with the $y$-axis, reflecting the high stiffness of the mount in this direction due to its elongated geometry and the direct constraint at the fixed base. The first modes exhibiting significant effective modal mass for translational displacement along the $y$-axis occur at frequencies higher than \SI{10}{\kilo\hertz}: \SI{11.4}{\kilo\hertz} for mode number $n=10$ and \SI{11.8}{\kilo\hertz}  for $n=11$. As a result, the resonance peaks observed within the \SIrange{0}{10}{\kilo\hertz} range for enforced acceleration along the $y$-axis are one to two orders of magnitude smaller than most resonances observed under $x$- and $z$-excitation.

\subsection{Passive stability of a fibre-tip resonator}
To evaluate the passive stability of the fibre-tip Fabry--Pérot resonator, we monitor its transmission under conditions such that the signal is directly proportional to cavity length. This is realised by operating the resonator in a low-finesse regime using \SI{1064}{\nano\meter} light, deliberately detuned from the mirror coating's design wavelength of \SI{780}{\nano\meter}. The lower reflectivity results in broader resonance profiles as shown in \autoref{fabryperot}. Before the transmission is recorded, the cavity length is adjusted with a piezoelectric actuator such that the transmitted optical power is at roughly \SI{50}{\percent} of its maximum value. At this point on the Airy distribution, the transmitted power varies approximately linearly with cavity length. A time trace of the transmission signal\footnote{It should be noted that this method is inherently sensitive to laser power fluctuations, since any change in the incident optical power is indistinguishable from an effective change in cavity length.} is recorded on a photodiode and the corresponding power spectral density is shown in \autoref{PSD}.
\begin{figure}[!t]
	%\vspace{0.7em}
	\centering
	\begin{tikzpicture}
	
	
	
	\begin{groupplot}[group style={group size=1 by 1, vertical sep=1cm, y descriptions at=edge right,},xmin=0,ymin=0,height=5.5cm,width=12cm,no markers,axis on top]
	
	\nextgroupplot
	[
	ymode=log,
	ymin=7e-1,
	ymax=1e4,
	ylabel style=markgrijs,
	yticklabel style=markgrijs,
	%ymin=1e-2, ymax=1e0,
	xlabel={Frequency {[\si{\kilo\hertz}]}},
	ylabel={Max. accel. {[\si{\milli\meter\per\second\squared}]}},
	xmin=0, xmax=10,
	ytick pos=right,
	%ymin=-0.1, ymax=1.1,
	%ymin=0, ymax=120,
	%xtick={0,20,40,60,80,100},
	%ytick={0,20,40,60,80,100,120},
	%legend pos=north west,
	ymajorgrids=false,
	xmajorgrids=false,
	grid style=dashed,
	legend cell align={left},
	legend style={
		at={(1, 1)},  % Adjust (x, y) to move it
		anchor=north east,  % Attach the legend to its top-right corner
		draw=none,           % Optional: Remove border around the legend
		fill=none,
		align=left
	}
	]
	
	\draw[color=markgrijs, dashed] 
	(axis cs:0.267, 7e-1) -- (axis cs:0.267, 40000);
	\draw[color=markgrijs, dashed] 
	(axis cs:1.105, 7e-1) -- (axis cs:1.105, 40000);
	\draw[color=markgrijs, dashed] 
	(axis cs:1.254, 7e-1) -- (axis cs:1.254, 40000);
	\draw[color=markgrijs, dashed] 
	(axis cs:2.940, 7e-1) -- (axis cs:2.940, 40000);
	\draw[color=markgrijs, dashed] 
	(axis cs:3.024, 7e-1) -- (axis cs:3.024, 40000);
	\draw[color=markgrijs, dashed] 
	(axis cs:3.893, 7e-1) -- (axis cs:3.893, 40000);
	\draw[color=markgrijs, dashed] 
	(axis cs:6.643, 7e-1) -- (axis cs:6.643, 40000);
	\draw[color=markgrijs, dashed] 
	(axis cs:6.935, 7e-1) -- (axis cs:6.935, 40000);
	
	
	
	\pgfplotsset{
		after end axis/.append code={
			\node[align=center] at (axis cs: 0.267, 20000) {\footnotesize\sc a};
			\node[align=center] at (axis cs: 1.105-0.06, 20000) {\footnotesize\sc b};
			\node[align=center] at (axis cs: 1.254+0.06, 20000) {\footnotesize\sc c};
			\node[align=center] at (axis cs: 2.940-0.1, 20000) {\footnotesize\sc d};
			\node[align=center] at (axis cs: 3.024+0.1, 20000) {\footnotesize\sc e};
			\node[align=center] at (axis cs: 3.893, 20000) {\footnotesize\sc f};
			\node[align=center] at (axis cs: 6.643, 20000) {\footnotesize\sc g};
			\node[align=center] at (axis cs: 6.935, 20000) {\footnotesize\sc h};
		}
	}
	
	
	\addplot[
	color=markgrijs,
	name path=xpointone,
	] table [col sep = comma, x expr = \thisrowno{0}/1000, y index = 1]{M_x_01.csv};
	
	\addplot+[draw=none,domain=0:10,name path=xzero] {1e-11}; 
	
	\addplot[color=markgrijs, opacity=0.25] fill between[
	of = xzero and xpointone,
	]; 
	
	\addplot[
	color=marklichtgrijs2!70!black,
	name path=zpointone,
	] table [col sep = comma, x expr = \thisrowno{0}/1000, y index = 1]{M_z_01.csv};
	
	\addplot+[draw=none,domain=0:10,name path=zzero] {1e-11}; 
	
	\addplot[color=markgrijs!10!white, opacity=0.8] fill between[
	of = zzero and zpointone,
	]; 
	
	\addplot[
	color=black,
	name path=ypointone,
	] table [col sep = comma, x expr = \thisrowno{0}/1000, y index = 1]{M_y_01.csv};
	
	\addplot+[draw=none,domain=0:10,name path=yzero] {1e-11}; 
	
	\addplot[color=markgrijs, opacity=0.5] fill between[
	of = yzero and ypointone,
	]; 
	
	\end{groupplot}
	
	\begin{groupplot}[group style={group size=1 by 1, vertical sep=1cm,y descriptions at=edge left},xmin=0,ymin=0,height=5.5cm,width=12cm,no markers]
	
	\nextgroupplot
	[
	ymode=log,
	ymin=1e-10,
	ymax=1e-4,
	ylabel style=markrood,
	yticklabel style=markrood,
	every y tick/.style={markrood},
	%ymin=1e-2, ymax=1e0,
	%xlabel={Frequency (\si{\kilo\hertz})},
	ylabel={Power spectral density {[\si{\volt\squared\per\hertz}]}},
	xmin=0, xmax=10,
	%ymin=-0.1, ymax=1.1,
	%ymin=0, ymax=120,
	%xtick={0,20,40,60,80,100},
	%ytick={0,20,40,60,80,100,120},
	%legend pos=north west,
	xtick=\empty,
	axis line style=transparent,
	ytick pos=left,
	ymajorgrids=true,
	xmajorgrids=false,
	grid style=dashed,
	y grid style={color=markrood, opacity=0.3},
	legend cell align={left},
	legend style={
		at={(1, 1)},  % Adjust (x, y) to move it
		anchor=north east,  % Attach the legend to its top-right corner
		draw=none,           % Optional: Remove border around the legend
		fill=none,
		align=left
	}
	]
	
	
	\addplot[
	color=markrood,
	] table [col sep = comma, x expr = \thisrowno{0}/1000, y index = 1]{PSD.csv};

	\node[align=center] at (rel axis cs: 0.335,0.11) {\footnotesize\sc \textcolor{marklichtgrijs2!70!black}{$z$}};
	\node[align=center] at (rel axis cs: 0.64,0.09) {\footnotesize\sc \textcolor{black}{$y$}};
	\node[align=center] at (rel axis cs: 0.42,0.13) {\footnotesize\sc \textcolor{markgrijs}{$x$}};
	
	\end{groupplot}
	
	\end{tikzpicture}
	\caption[Power spectral density of the fibre-tip resonator transmission signal at \SI{1064}{\nano\meter}]{Power spectral density of the fibre-tip resonator transmission signal at \SI{1064}{\nano\meter}. The grey reference traces correspond to the simulated modal frequency response curves ($x$, $y$ and $z$) as shown in \autoref{freqresponseMarw}. Labels are references to the modes as indicated in \autoref{freqresponseMarw} and \autoref{marwanmodes}.}
	\label{PSD} 
\end{figure}
As a reference, the modal frequency response traces from \autoref{freqresponseMarw} are shown in grey. The power spectral density of the transmission signal reveals that, under ambient conditions, environmental vibrational noise perturbs the cavity length most strongly at frequencies near the expected modal resonances.

It is important to note that the transmission spectrum in \autoref{PSD} is not directly comparable to a modal frequency response curve. The measured signal reflects the product of the environmental vibration spectrum with the transfer function of the entire vacuum chamber and mounting assembly. Therefore, we do not expect an exact correspondence with the simulated modal frequency response. Nonetheless, we can identify a number of common features. The most prominent peak in the transmission spectrum appears at \SI{200}{\hertz}, which lies reasonably close to the first normal mode of the mounting structure at \SI{267}{\hertz} \textsc{(a)}. The discrepancy can be attributed to the additional mass of a small steel fibre holder that is fastened to the free end of the mounting structure with four screws (not shown in \autoref{marwanmodes}). When material is added at the free end of a cantilever, two competing effects arise: the added inertia lowers the resonant frequency, while the increased local stiffness raises it. For a local thickening near the free end of a cantilever, the inertial effect dominates, resulting a net downshift of the fundamental bending frequency. A second pronounced feature in the transmission spectrum is visible at \SI{1220}{\hertz}, closely matching modes \textsc{(b)} and \textsc{(c)} at \SI{1105}{\hertz} and \SI{1254}{\hertz}, respectively. At higher frequencies, two peaks are observed at \SI{3330}{\hertz} and \SI{3510}{\hertz}, somewhat above the simulated pair of resonances at \SI{2940}{\hertz} \textsc{(d)} and \SI{3024}{\hertz} \textsc{(e)}. This upward shift can be explained by the stiffening effect of the fibre support piece, which is more significant for these bending modes. In contrast, modes \textsc{(g)} and \textsc{(h)} are only weakly visible in the transmission spectrum, most likely because the environmental background provides limited excitation in this frequency range.

\subsection[Actuator transfer function of a fibre-tip resonator mount]{Actuator transfer function of a fibre-tip\hspace{4em} resonator mount}
In the previous section, the passive stability of the fibre-tip Fabry--Pérot resonator was characterised by analysing the transmission signal under ambient conditions. This approach, however, is inherently limited by the non-uniform power spectrum of the environmental noise. To obtain a more controlled and quantitative characterisation of the mechanical response, we measure the actuator transfer function of the fibre-tip resonator mount. For this purpose, one of the shear piezoelectric actuators that holds one end of the fibre-tip resonator is used to drive\footnote{Within the vacuum chamber, this configuration provides the only available means to actuate the mount mechanically.} the cavity and its mount sinusoidally, while the drive frequency is swept up to \SI{10}{\kilo\hertz} with an arbitrary waveform generator (AWG). The sinusoidal drive induces a corresponding oscillation of the cavity length and the response amplitude is expected to increase when the drive frequency coincides with a mechanical resonance of the mount. Strictly speaking, this measurement also probes resonances of the piezoelectric actuator itself, but these occur far above the frequency range of interest. The specified unloaded resonance frequency of the Noliac CSAP02 shear piezoelectric actuator is \SI{1.75}{\mega\hertz}. While the additional mass of the fibres and holders bonded to the actuator may reduce this resonance frequency, it is not expected to shift into the sub-\SI{10}{\kilo\hertz} range explored here. As in the passive stability analysis, the resonator is illuminated at \SI{1064}{\nano\meter} and a second piezoelectric actuator is used to stabilise the cavity length at the half-transmission point of a resonance slope, ensuring that the transmitted optical power remains approximately linear with respect to cavity length. The transmitted signal from the cavity is recorded on a photodiode and analysed with a spectrum analyser. By comparing the amplitude and phase of the photodiode signal to those of the AWG reference, we directly obtain the magnitude and phase of the transfer function $H(i\cdot 2\pi\nu)$, defined as the ratio of the cavity transmission response (photodiode output voltage) to the applied drive voltage before amplification. The results are shown in \autoref{transferfunction}.
\begin{figure}[!t]
	%\vspace{0.7em}
	\centering
	\begin{tikzpicture}
	
	
	
	\begin{groupplot}[group style={group size=1 by 2, vertical sep=0.5cm, y descriptions at=edge right,},xmin=0,ymin=0,height=4.5cm,width=12cm,no markers,axis on top]
	
	\nextgroupplot
	[
	ymode=log,
	ymin=7e-1,
	ymax=1e4,
	ylabel style={markgrijs, xshift=-4em},
	yticklabel style=markgrijs,
	%ymin=1e-2, ymax=1e0,
	%xlabel={Frequency (\si{\kilo\hertz})},
	ylabel={Max. accel. {[\si{\milli\meter\per\second\squared}]}},
	xmin=0, xmax=10,
	ytick pos=right,
	xticklabels={,,},
	%ymin=-0.1, ymax=1.1,
	%ymin=0, ymax=120,
	%xtick={0,20,40,60,80,100},
	%ytick={0,20,40,60,80,100,120},
	%legend pos=north west,
	ymajorgrids=false,
	xmajorgrids=false,
	grid style=dashed,
	legend cell align={left},
	legend style={
		at={(1, 1)},  % Adjust (x, y) to move it
		anchor=north east,  % Attach the legend to its top-right corner
		draw=none,           % Optional: Remove border around the legend
		fill=none,
		align=left,
		%y label style={at={(axis description cs: -0.12, 0.5)},anchor=base}
	}
	]
	
	\draw[color=markgrijs, dashed] 
	(axis cs:0.267, 7e-1) -- (axis cs:0.267, 40000);
	\draw[color=markgrijs, dashed] 
	(axis cs:1.105, 7e-1) -- (axis cs:1.105, 40000);
	\draw[color=markgrijs, dashed] 
	(axis cs:1.254, 7e-1) -- (axis cs:1.254, 40000);
	\draw[color=markgrijs, dashed] 
	(axis cs:2.940, 7e-1) -- (axis cs:2.940, 40000);
	\draw[color=markgrijs, dashed] 
	(axis cs:3.024, 7e-1) -- (axis cs:3.024, 40000);
	\draw[color=markgrijs, dashed] 
	(axis cs:3.893, 7e-1) -- (axis cs:3.893, 40000);
	\draw[color=markgrijs, dashed] 
	(axis cs:6.643, 7e-1) -- (axis cs:6.643, 40000);
	\draw[color=markgrijs, dashed] 
	(axis cs:6.935, 7e-1) -- (axis cs:6.935, 40000);
	
	
	\addplot[
	color=markgrijs,
	name path=xpointone,
	] table [col sep = comma, x expr = \thisrowno{0}/1000, y index = 1]{M_x_01.csv};
	
	\addplot+[draw=none,domain=0:10,name path=xzero] {1e-11}; 
	
	\addplot[color=markgrijs, opacity=0.25] fill between[
	of = xzero and xpointone,
	]; 
	
	\addplot[
	color=marklichtgrijs2!70!black,
	name path=zpointone,
	] table [col sep = comma, x expr = \thisrowno{0}/1000, y index = 1]{M_z_01.csv};
	
	\addplot+[draw=none,domain=0:10,name path=zzero] {1e-11}; 
	
	\addplot[color=markgrijs!10!white, opacity=0.8] fill between[
	of = zzero and zpointone,
	]; 
	
	\addplot[
	color=black,
	name path=ypointone,
	] table [col sep = comma, x expr = \thisrowno{0}/1000, y index = 1]{M_y_01.csv};
	
	\addplot+[draw=none,domain=0:10,name path=yzero] {1e-11}; 
	
	\addplot[color=markgrijs, opacity=0.5] fill between[
	of = yzero and ypointone,
	]; 
	
	\nextgroupplot
	[
	ymode=log,
	ymin=7e-1,
	ymax=1e4,
	ylabel style=markgrijs,
	yticklabel style=markgrijs,
	%ymin=1e-2, ymax=1e0,
	xlabel={Frequency $\nu$ {[\si{\kilo\hertz}]}},
	%ylabel={Max. accel. (\si{\milli\meter\per\second\squared})},
	xmin=0, xmax=10,
	ytick pos=right,
	%ymin=-0.1, ymax=1.1,
	%ymin=0, ymax=120,
	%xtick={0,20,40,60,80,100},
	%ytick={0,20,40,60,80,100,120},
	%legend pos=north west,
	ymajorgrids=false,
	xmajorgrids=false,
	grid style=dashed,
	legend cell align={left},
	legend style={
		at={(1, 1)},  % Adjust (x, y) to move it
		anchor=north east,  % Attach the legend to its top-right corner
		draw=none,           % Optional: Remove border around the legend
		fill=none,
		align=left
	},
	y label style={at={(axis description cs: -0.12, 0.5)},anchor=base}
	]
	
	\draw[color=markgrijs, dashed] 
	(axis cs:0.267, 7e-1) -- (axis cs:0.267, 40000);
	\draw[color=markgrijs, dashed] 
	(axis cs:1.105, 7e-1) -- (axis cs:1.105, 40000);
	\draw[color=markgrijs, dashed] 
	(axis cs:1.254, 7e-1) -- (axis cs:1.254, 40000);
	\draw[color=markgrijs, dashed] 
	(axis cs:2.940, 7e-1) -- (axis cs:2.940, 40000);
	\draw[color=markgrijs, dashed] 
	(axis cs:3.024, 7e-1) -- (axis cs:3.024, 40000);
	\draw[color=markgrijs, dashed] 
	(axis cs:3.893, 7e-1) -- (axis cs:3.893, 40000);
	\draw[color=markgrijs, dashed] 
	(axis cs:6.643, 7e-1) -- (axis cs:6.643, 40000);
	\draw[color=markgrijs, dashed] 
	(axis cs:6.935, 7e-1) -- (axis cs:6.935, 40000);
	
	
	
	\pgfplotsset{
		after end axis/.append code={
			\node[align=center] at (axis cs: 0.267, 20000) {\footnotesize\sc a};
			\node[align=center] at (axis cs: 1.105-0.06, 20000) {\footnotesize\sc b};
			\node[align=center] at (axis cs: 1.254+0.06, 20000) {\footnotesize\sc c};
			\node[align=center] at (axis cs: 2.940-0.1, 20000) {\footnotesize\sc d};
			\node[align=center] at (axis cs: 3.024+0.1, 20000) {\footnotesize\sc e};
			\node[align=center] at (axis cs: 3.893, 20000) {\footnotesize\sc f};
			\node[align=center] at (axis cs: 6.643, 20000) {\footnotesize\sc g};
			\node[align=center] at (axis cs: 6.935, 20000) {\footnotesize\sc h};
		}
	}
	
	
	
	\addplot[
	color=markgrijs,
	name path=xpointone,
	] table [col sep = comma, x expr = \thisrowno{0}/1000, y index = 1]{M_x_01.csv};
	
	\addplot+[draw=none,domain=0:10,name path=xzero] {1e-11}; 
	
	\addplot[color=markgrijs, opacity=0.25] fill between[
	of = xzero and xpointone,
	]; 
	
	\addplot[
	color=marklichtgrijs2!70!black,
	name path=zpointone,
	] table [col sep = comma, x expr = \thisrowno{0}/1000, y index = 1]{M_z_01.csv};
	
	\addplot+[draw=none,domain=0:10,name path=zzero] {1e-11}; 
	
	\addplot[color=markgrijs!10!white, opacity=0.8] fill between[
	of = zzero and zpointone,
	]; 
	
	\addplot[
	color=black,
	name path=ypointone,
	] table [col sep = comma, x expr = \thisrowno{0}/1000, y index = 1]{M_y_01.csv};
	
	\addplot+[draw=none,domain=0:10,name path=yzero] {1e-11}; 
	
	\addplot[color=markgrijs, opacity=0.5] fill between[
	of = yzero and ypointone,
	]; 
	
	\end{groupplot}
	
	\begin{groupplot}[group style={group size=1 by 2, vertical sep=0.5cm,y descriptions at=edge left},xmin=0,ymin=0,height=4.5cm,width=12cm,no markers]
	
	\nextgroupplot
	[
	%ymode=log,
	%ymin=1e-10,
	%ymax=1e-4,
	ylabel style=markrood,
	yticklabel style=markrood,
	every y tick/.style={markrood},
	%ymin=1e-2, ymax=1e0,
	%xlabel={Frequency (\si{\kilo\hertz})},
	ylabel={$\vert H(i\cdot 2\pi\nu)\vert$ {[\si{\deci B}]}},
	xmin=0, xmax=10,
	%ytick={0,0.05,0.1},
	%yticklabels = {0,0.05,0.1},
	ymin=-37, ymax=-18,
	%ymin=0, ymax=120,
	%xtick={0,20,40,60,80,100},
	%ytick={0,20,40,60,80,100,120},
	%legend pos=north west,
	xtick=\empty,
	axis line style=transparent,
	ytick pos=left,
	ymajorgrids=true,
	xmajorgrids=false,
	grid style=dashed,
	y grid style={color=markrood, opacity=0.3},
	legend cell align={left},
	legend style={
		at={(1, 1)},  % Adjust (x, y) to move it
		anchor=north east,  % Attach the legend to its top-right corner
		draw=none,           % Optional: Remove border around the legend
		fill=none,
		align=left
	},
	y label style={at={(axis description cs: -0.12, 0.5)},anchor=base}
	]
	
	
	\addplot[
	color=markrood,
	] table [col sep = comma, x expr = \thisrowno{0}/1000, y expr={20 * ln(\thisrowno{1}) / ln(10)}]{TRANSFERFUNCTION.csv};
	
	\node[align=center] at (rel axis cs: 0.32,0.15) {\footnotesize\sc \textcolor{marklichtgrijs2!80!black}{$z$}};
	\node[align=center] at (rel axis cs: 0.98,0.10) {\footnotesize\sc \textcolor{black}{$y$}};
	\node[align=center] at (rel axis cs: 0.41,0.15) {\footnotesize\sc \textcolor{markgrijs}{$x$}};
	
	\nextgroupplot
	[
	%ymode=log,
	%ymin=1e-10,
	%ymax=1e-4,
	ylabel style=markdiepblauw,
	yticklabel style=markdiepblauw,
	every y tick/.style={markdiepblauw},
	%ymin=1e-2, ymax=1e0,
	%xlabel={Frequency (\si{\kilo\hertz})},
	ylabel={$\arg{(H(i\cdot 2\pi\nu))}$ {[\si{\degree}]}},
	xmin=0, xmax=10,
	ymin=-150, ymax=0,
	%ymin=0, ymax=120,
	%xtick={0,20,40,60,80,100},
	%ytick={0,20,40,60,80,100,120},
	%legend pos=north west,
	xtick=\empty,
	axis line style=transparent,
	ytick pos=left,
	ymajorgrids=true,
	xmajorgrids=false,
	grid style=dashed,
	y grid style={color=markdiepblauw, opacity=0.3},
	legend cell align={left},
	legend style={
		at={(1, 1)},  % Adjust (x, y) to move it
		anchor=north east,  % Attach the legend to its top-right corner
		draw=none,           % Optional: Remove border around the legend
		fill=none,
		align=left
	},
	y label style={at={(axis description cs: -0.12, 0.5)},anchor=base}
	]
	
	
	\addplot[
	color=markdiepblauw,
	] table [col sep = comma, x expr = \thisrowno{0}/1000, y expr={\thisrowno{2}-180}]{TRANSFERFUNCTION.csv};
	
	\node[align=center] at (rel axis cs: 0.32,0.15) {\footnotesize\sc \textcolor{marklichtgrijs2!70!black}{$z$}};
	\node[align=center] at (rel axis cs: 0.98,0.10) {\footnotesize\sc \textcolor{black}{$y$}};
	\node[align=center] at (rel axis cs: 0.41,0.15) {\footnotesize\sc \textcolor{markgrijs}{$x$}};
	
	\end{groupplot}
	
	\end{tikzpicture}
	\caption[Magnitude and phase of the transfer function that corresponds to the transmission at \SI{1064}{\nano\meter} of a piezo-driven fibre-tip resonator]{Magnitude \textsc{(top)} and phase \textsc{(bottom)} of the transfer function $H(i\cdot 2\pi\nu)$ that corresponds to the transmission signal at \SI{1064}{\nano\meter} of a fibre-tip resonator that is driven sinusoidally with a single sheer piezoelectric actuator. The grey reference traces correspond to the simulated modal frequency response curves for enforced acceleration along the axes $x$, $y$ (darkest) and $z$ (lightest) as shown in \autoref{freqresponseMarw}. Labels are references to the modes as indicated in \autoref{freqresponseMarw} and \autoref{marwanmodes}.}
	\label{transferfunction} 
\end{figure}
The measured transfer function of the fibre-tip resonator mount exhibits sharp resonance-antiresonance pairs, each one of which qualitatively resembles the characteristic response of a two-mode mechanical system with collocated actuation and sensing (see \autoref{fig:tf_groupplot} in \autoref{app:modalanalysis}). In the present case, the shear piezoelectric actuator applies a local force along the $y$-axis of the mount at a position within \SI{5}{\milli\meter} of the cavity centre. The cavity transmission signal, which does not directly correspond to a measurement of displacement, rather responds to bending of the mount at this location, which results in the actuation and sensing points being nearly collocated.

If we examine the transfer function, the most prominent features appear close to the predicted resonances from the frequency response analysis for enforced acceleration of the base along the $y$-axis, where the visible modes associated with this direction of actuation are \textsc{(c)}, \textsc{(e)} and \textsc{(g)}. This behaviour is expected, since the shear piezoelectric actuator applies a harmonic force along the $y$-axis, though at a position near the free end of the cantilever. By contrast, the two lowest-frequency resonances that dominate the passive stability spectrum in \autoref{PSD} leave almost no signature in the transfer function, as they are primarily excited via enforced acceleration along the $z$- and $x$-axes, respectively (see \autoref{freqresponseMarw}). These resonances are thus excited more effectively by environmental vibrations that couple through the base in multiple directions. Interestingly, even though most of the other resonances \textsc{(a, b, d, f)} remain essentially invisible in the transfer function, we do observe mode \textsc{(h)} alongside mode \textsc{(g)}. In the simulations, however, mode \textsc{(h)} only appears for enforced acceleration along the $x$-axis. This suggests that a small component of the actuator force may couple into the $x$-direction, possibly due to slight misalignment of the piezoelectric element.

Finally, additional spectral features appear around \SI{9}{\kilo\hertz}. Since these do not emerge in the simulated modal response curves, they are most likely associated with the fibre-holder structure, which was not included in the finite-element model. These observations demonstrate that actuator transfer function measurements provide complementary information to the passive stability analysis. The latter reveals which modes are most strongly excited by ambient disturbances, while the former identifies which modes are efficiently driven and sensed in the actuator-cavity configuration. Taken together, these perspectives yield a more comprehensive picture of the dynamic behaviour of the fibre-tip resonator mount.

\subsection{Design of a glass-substrate resonator mount}
\label{subsec:mountingoverview}
The fibre-tip resonator mount, while simple in its design, suffers from poor mechanical stability. A dense set of low-frequency resonances spread across the simulated frequency response and measured transfer function makes the cavity highly sensitive to environmental vibrations and difficult to operate reliably. In particular, the presence of strong resonances complicates active feedback stabilisation of cavity length, since the controller gain must be reduced around these resonances to maintain stability, which in turn limits the achievable feedback bandwidth. To overcome these limitations, we developed a new mount for a glass-substrate resonator that emphasises passive stability by maximising the steel mass inside the glass cell and by shaping the structure to break degeneracies in the mechanical response spectrum. This design philosophy aims to shift resonances out of the critical low-frequency regime and provide a more robust platform for stabilising cavity length.

The mounting structure, shown in \autoref{fig:fingeroverview}, is machined from 316LN stainless steel and houses the cavity module (see \autoref{cavitymodulesection}), which slides in through the central square aperture and is fixed at the correct position with six countersunk screws. In addition, the mounting structure hosts a compact pyramidal mirror assembly (see \autoref{sec:pyramidal_mot}), held in place by a custom aluminium holder inserted from the side.
\begin{figure}[t]
	\centering
	\begin{tikzpicture}[remember picture]
	\node[anchor=south west,inner sep=0] (image1) at (0,0) {\includegraphics[width=1\textwidth]{{"finger"}.pdf}};
	
	% left four
	\begin{scope}[x={(image1.south east)},y={(image1.north west)}]
	% draw a grid
	%\draw[help lines,xstep=.1,ystep=.1,overlay] (0,0) grid (1,1);
	% draw ticks
	%\foreach \x in {0,1,...,9} { \node [anchor=north,overlay] at (\x/10,0) {0.\x}; }
	%\foreach \y in {0,1,...,9} { \node [anchor=east,overlay] at (0,\y/10) {0.\y}; }
	
	% labels
	\node[align=center, rotate=0] at (0.05,0.9) {\footnotesize\sc (a)};
	\node[align=center, rotate=0] at (0.05,0.42) {\footnotesize\sc (b)};
	\node[align=center, rotate=0] at (0.67,0.42) {\footnotesize\sc (c)};
	
	\end{scope}
	
	\end{tikzpicture}
	\caption[Cavity mounting structure and integration into the vacuum system]{Cavity mounting structure and integration into the vacuum system. \textsc{(a)} The cavity mounting structure is designed to secure the cavity module with six screws. A small mount inserted from the side holds four mirrors in a pyramidal configuration to form a MOT from which to load the dipole trap. The insets on the right illustrate the assembly of the pyramidal MOT mirrors within the mount. The mounting structure features multiple openings to provide optical access for (horizontal) cooling and (vertical) trapping beams. \textsc{(b)} Integration of the cavity mounting structure into the vacuum system. A cuboid glass cell encloses the mounting structure, which is bolted to a flange attached to the vacuum chamber. This flange also secures the mounting structure with six smaller bolts (see inset). \textsc{(c)} A partial view of the mounting structure inside the glass cell, showing a pair of high-NA lenses that focus and recollimate the trapping beam (shown in red).}
	\label{fig:fingeroverview} 
\end{figure}
Multiple apertures are machined into the mount to allow optical access for horizontal cooling beams, the vertically propagating trapping beam and cavity coupling beams. The entire structure is enclosed in a cuboid glass cell equipped with an anti-reflective coating as discussed in \autoref{setupoverview}. Both the mounting structure and glass cell are secured to a custom vacuum flange with two sets of six bolts.

\subsection{Modal frequency response analysis of a glass-substrate resonator mount}
The modal frequency response of the glass-substrate resonator mount (including cavity module) is computed with the same finite-element analysis settings used for the fibre-tip resonator mount (enforced base acceleration of \SI{1}{\milli\metre\per\second\squared} and identical damping cases), but on a mesh of \num{330715} elements to maintain element size. The effective modal mass ratios are reported in \autoref{tab:modepercentFinalfinger} (see \autoref{app:effectivemodalmass}), while the frequency response curves and the mode shapes below \SI{10}{\kilo\hertz} are shown in \autoref{freqresponseFF} and \autoref{FFmodes},
\begin{figure}[!t]
	%\vspace{0.7em}
	\centering
	\begin{tikzpicture}
	
	\begin{groupplot}[group style={group size=1 by 3, vertical sep=0.5cm},xmin=0,ymin=0,height=4cm,width=13.5cm,no markers]
	
	\iffalse
	\nextgroupplot
	[
	ylabel={},
	xmin=0, xmax=1,
	ymin=0, ymax=1,
	xticklabels={,\textcolor{white}{0},},
	yticklabels={,,},
	%legend pos=north west,
	ymajorgrids=true,
	xmajorgrids=true,
	grid style=dashed,
	%axis lines=none,        % No axis lines
	%tick style=none,        % No ticks
	%xtick=\empty,           % No x-axis ticks
	%ytick=\empty,           % No y-axis ticks
	xlabel={},              % No x-axis label
	ylabel={},              % No y-axis label
	%enlargelimits=true,     % Keeps the plot within the bounds
	]
	
	\addplot graphics [
	xmin=0, xmax=1,  % Stretch to fit the x-axis limits
	ymin=0, ymax=1    % Stretch to fit the y-axis limits
	] {FINALFINGER.png};
	\fi
	
	
	\nextgroupplot
	[
	ymode=log,
	ymin=1e-0,
	ymax=1e4,
	%xlabel={Frequency},
	%ylabel={Max. linear acceleration},
	xmin=0, xmax=10,
	xticklabels={,,},
	%ymin=-0.1, ymax=1.1,
	%ymin=0, ymax=120,
	%xtick={0,20,40,60,80,100},
	%ytick={0,20,40,60,80,100,120},
	%legend pos=north west,
	ymajorgrids=true,
	xmajorgrids=true,
	grid style=dashed,
	legend cell align={left},
	legend style={
		at={(1, 1)},  % Adjust (x, y) to move it
		anchor=north east,  % Attach the legend to its top-right corner
		draw=none,           % Optional: Remove border around the legend
		fill=none,
		align=left
	}
	]
	
	%\node[align=center] at (axis cs: 1.890, 3500) {\footnotesize A};
	\node[align=center] at (axis cs: 1.979, 3500) {\footnotesize\sc b};
	\node[align=center] at (axis cs: 5.893, 2500) {\footnotesize\sc c};
	\node[align=center] at (axis cs: 6.337, 300) {\footnotesize\sc d};
	\node[align=center] at (axis cs: 6.466, 700) {\footnotesize\sc e};
	\node[align=center] at (axis cs: 8.530+0.2, 90) {\footnotesize\sc f};
	
	
	\addplot[
	color=markrood,
	name path=xpointone,
	] table [col sep = comma, x expr = \thisrowno{0}/1000, y index = 1]{F_x_01.csv};
	
	\addplot[
	color=rainbow2of8,
	name path=xone,
	] table [col sep = comma, x expr = \thisrowno{0}/1000, y index = 1]{F_x_1.csv};
	
	\addplot[
	color=rainbow4of8,
	name path=xfive,
	] table [col sep = comma, x expr = \thisrowno{0}/1000, y index = 1]{F_x_5.csv};
	
	\addplot+[draw=none,domain=0:10,name path=xzero] {0.01}; 
	
	\addplot[color=markrood, opacity=0.5] fill between[
	of = xone and xpointone,
	]; 
	
	\addplot[color=rainbow2of8, opacity=0.5] fill between[
	of = xfive and xone,
	]; 
	
	\addplot[color=rainbow4of8, opacity=0.5] fill between[
	of = xzero and xfive,
	]; 
	
	\node at (rel axis cs: 0.39, 0.55) [anchor=center] {\includegraphics[width=3.1cm]{FINALFINGERtrans.png}};
	
	\addlegendentry{\footnotesize \hspace{0.3em} $G=\SI{0.1}{\percent}$}
	\addlegendentry{\footnotesize \hspace{0.3em} $G=\SI{1}{\percent}$}
	
	
	\draw[<->] (rel axis cs: 0.54, 0.75) --++ (150:0.5cm);
	\node[align=left] at (rel axis cs: 0.532, 0.88) {\footnotesize $x$};
	
	\nextgroupplot
	[
	ymode=log,
	ymin=1e-0,
	ymax=1e2,
	%ymin=1e-2, ymax=1e0,
	%xlabel={Frequency},
	ylabel={Maximum linear acceleration magnitude {[\si{\milli\meter\per\second\squared}]}},
	xmin=0, xmax=10,
	xticklabels={,,},
	%ymin=-0.1, ymax=1.1,
	%ymin=0, ymax=120,
	%xtick={0,20,40,60,80,100},
	%ytick={0,20,40,60,80,100,120},
	%legend pos=north west,
	ymajorgrids=true,
	xmajorgrids=true,
	grid style=dashed,
	legend cell align={left},
	legend style={
		at={(1, 1)},  % Adjust (x, y) to move it
		anchor=north east,  % Attach the legend to its top-right corner
		draw=none,           % Optional: Remove border around the legend
		fill=none,
		align=left
	}
	]
	
	%\node[align=center] at (axis cs: 1.890, 3500) {\footnotesize A};
	\node[align=center] at (axis cs: 1.979, 6) {\footnotesize\sc b};
	\node[align=center] at (axis cs: 5.893, 20) {\footnotesize\sc c};
	\node[align=center] at (axis cs: 6.337, 61) {\footnotesize\sc d};
	\node[align=center] at (axis cs: 6.466, 42) {\footnotesize\sc e};
	\node[align=center] at (axis cs: 8.530, 13) {\footnotesize\sc f};
	
	\addplot[
	forget plot,
	color=markrood,
	name path=ypointone,
	] table [col sep = comma, x expr = \thisrowno{0}/1000, y index = 1]{F_y_01.csv};
	
	\addplot[
	forget plot,
	color=rainbow2of8,
	name path=yone,
	] table [col sep = comma, x expr = \thisrowno{0}/1000, y index = 1]{F_y_1.csv};
	
	\addplot[
	color=rainbow4of8,
	name path=yfive,
	] table [col sep = comma, x expr = \thisrowno{0}/1000, y index = 1]{F_y_5.csv};
	
	\addplot+[draw=none,domain=0:10,name path=yzero] {0.01}; 
	
	\addplot[color=markrood, opacity=0.5] fill between[
	of = yone and ypointone,
	]; 
	
	\addplot[color=rainbow2of8, opacity=0.5] fill between[
	of = yfive and yone,
	]; 
	
	\addplot[color=rainbow4of8, opacity=0.5] fill between[
	of = yzero and yfive,
	]; 
	
	\node at (rel axis cs: 0.39, 0.55) [anchor=center] {\includegraphics[width=3.1cm]{FINALFINGERtrans.png}};
	
	%\addlegendentry{\small \hspace{0.3em} $G=\SI{0.1}{\percent}$}
	%\addlegendentry{\small \hspace{0.3em} $G=\SI{1}{\percent}$}
	%\addlegendentry{\small \hspace{0.3em} $G=\SI{5}{\percent}$}
	
	\addlegendentry{\footnotesize \hspace{0.3em} $G=\SI{5}{\percent}$}
	
	\draw[<->] (rel axis cs: 0.505, 0.76) --++ (30:0.5cm);
	\node[align=left] at (rel axis cs: 0.515, 0.895) {\footnotesize $y$};
	
	\nextgroupplot
	[
	ymode=log,
	ymin=1e-0,
	ymax=1e4,
	%ymin=1e-2, ymax=1e0,
	xlabel={Enforced acceleration frequency {[\si{\kilo\hertz}]}},
	%ylabel={Max. linear acceleration},
	xmin=0, xmax=10,
	%ymin=-0.1, ymax=1.1,
	%ymin=0, ymax=120,
	%xtick={0,20,40,60,80,100},
	%ytick={0,20,40,60,80,100,120},
	%legend pos=north west,
	ymajorgrids=true,
	xmajorgrids=true,
	grid style=dashed,
	legend cell align={left},
	legend style={
		at={(1, 1)},  % Adjust (x, y) to move it
		anchor=north east,  % Attach the legend to its top-right corner
		draw=none,           % Optional: Remove border around the legend
		fill=none,
		align=left
	}
	]
	
	\node[align=center] at (axis cs: 1.890, 3500) {\footnotesize\sc a};
	%\node[align=center] at (axis cs: 1.979, 3500) {\footnotesize B};
	\node[align=center] at (axis cs: 5.893, 55) {\footnotesize\sc c};
	\node[align=center] at (axis cs: 6.337, 5000) {\footnotesize\sc d};
	\node[align=center] at (axis cs: 6.466, 2500) {\footnotesize\sc e};
	\node[align=center] at (axis cs: 8.530, 300) {\footnotesize\sc f};
	
	\addplot[
	color=markrood,
	name path=zpointone,
	] table [col sep = comma, x expr = \thisrowno{0}/1000, y index = 1]{F_z_01.csv};
	
	\addplot[
	color=rainbow2of8,
	name path=zone,
	] table [col sep = comma, x expr = \thisrowno{0}/1000, y index = 1]{F_z_1.csv};
	
	\addplot[
	color=rainbow4of8,
	name path=zfive,
	] table [col sep = comma, x expr = \thisrowno{0}/1000, y index = 1]{F_z_5.csv};
	
	\addplot+[draw=none,domain=0:10,name path=zzero] {0.01}; 
	
	\addplot[color=markrood, opacity=0.5] fill between[
	of = zone and zpointone,
	]; 
	
	\addplot[color=rainbow2of8, opacity=0.5] fill between[
	of = zfive and zone,
	]; 
	
	\addplot[color=rainbow4of8, opacity=0.5] fill between[
	of = zzero and zfive,
	]; 
	
	\node at (rel axis cs: 0.39, 0.55) [anchor=center] {\includegraphics[width=3.1cm]{FINALFINGERtrans.png}};
	
	\draw[<->] (rel axis cs: 0.515, 0.74) --++ (90:0.5cm);
	\node[align=left] at (rel axis cs: 0.535, 0.86) {\footnotesize $z$};
	
	%\addlegendentry{\small \hspace{0.3em} $G=\SI{0.1}{\percent}$}
	%\addlegendentry{\small \hspace{0.3em} $G=\SI{1}{\percent}$}
	%\addlegendentry{\small \hspace{0.3em} $G=\SI{5}{\percent}$}
	
	\end{groupplot}
	
	
	
	\end{tikzpicture}
	\caption[Modal frequency response analysis of the glass-substrate resonator mounting structure]{Modal frequency response analysis of the glass-substrate resonator mounting structure. The graphs show the maximum linear acceleration magnitude of the structure versus the frequency of enforced acceleration at the base in three different directions: $x$ \textsc{(top)}, $y$ \textsc{(middle)} and $z$ \textsc{(bottom)}. In each direction, traces are shown for three different values of the structural damping coefficient $G$. The magnitude of the enforced acceleration equals \SI{1}{\milli\meter\per\second\squared}}
	\label{freqresponseFF} 
\end{figure}
respectively.
\begin{figure}[!t]
	\vspace{0em}
	\centering
	\begin{subfigure}[b]{.45\linewidth}
		\begin{tikzpicture}
		\centering
		\node[inner sep=0] (image) at (0,0) {\includegraphics[width=5cm] {F1F.png}};;
		\begin{scope}[x={(image.south east)},y={(image.north west)}]
		\node at (0.2, 0.7) {\footnotesize \SI{1890}{\hertz}};
		\end{scope}
		\end{tikzpicture}
		\centering
		\caption{}
		\label{a}
	\end{subfigure}
	\begin{subfigure}[b]{.45\linewidth}
		\begin{tikzpicture}
		\centering
		\node[inner sep=0] (image) at (0,0) {\includegraphics[width=5cm] {F2F.png}};;
		\begin{scope}[x={(image.south east)},y={(image.north west)}]
		\node at (0.2, 0.7) {\footnotesize \SI{1979}{\hertz}};
		\end{scope}
		\end{tikzpicture}
		\centering
		\caption{}
		\label{b}
	\end{subfigure}
	\\
	\begin{subfigure}[b]{.45\linewidth}
		\begin{tikzpicture}
		\centering
		\node[inner sep=0] (image) at (0,0) {\includegraphics[width=5cm] {F3F.png}};;
		\begin{scope}[x={(image.south east)},y={(image.north west)}]
		\node at (0.2, 0.7) {\footnotesize \SI{5893}{\hertz}};
		\end{scope}
		\end{tikzpicture}
		\centering
		\caption{}
		\label{c}
	\end{subfigure}
	\begin{subfigure}[b]{.45\linewidth}
		\begin{tikzpicture}
		\centering
		\node[inner sep=0] (image) at (0,0) {\includegraphics[width=5cm] {F4F.png}};;
		\begin{scope}[x={(image.south east)},y={(image.north west)}]
		\node at (0.2, 0.7) {\footnotesize \SI{6337}{\hertz}};
		\end{scope}
		\end{tikzpicture}
		\centering
		\caption{}
		\label{c}
	\end{subfigure}
	\\
	\begin{subfigure}[b]{.45\linewidth}
		\begin{tikzpicture}
		\centering
		\node[inner sep=0] (image) at (0,0) {\includegraphics[width=5cm] {F5F.png}};;
		\begin{scope}[x={(image.south east)},y={(image.north west)}]
		\node at (0.2, 0.7) {\footnotesize \SI{6466}{\hertz}};
		\end{scope}
		\end{tikzpicture}
		\centering
		\caption{}
		\label{a}
	\end{subfigure}
	\begin{subfigure}[b]{.45\linewidth}
		\begin{tikzpicture}
		\centering
		\node[inner sep=0] (image) at (0,0) {\includegraphics[width=5cm] {F6F.png}};;
		\begin{scope}[x={(image.south east)},y={(image.north west)}]
		\node at (0.2, 0.7) {\footnotesize \SI{8530}{\hertz}};
		\draw[->] (1, 0.15) --++ (90:0.5cm) node[above, anchor=south] {\footnotesize $z$};
		\draw[->] (1, 0.15) --++ (150:0.5cm) node[above, anchor=base] {\footnotesize  $x$\,\, \phantom{x}};
		\draw[->] (1, 0.15) --++ (210:0.5cm) node[above, anchor=mid] {\footnotesize $y$\,\, \phantom{x}};
		\end{scope}
		\end{tikzpicture}
		\centering
		\caption{}
		\label{b}
	\end{subfigure}
	\caption[Normal modes below \SI{10}{\kilo\hertz} of the glass-substrate resonator mounting structure]{Normal modes and corresponding frequencies below \SI{10}{\kilo\hertz} of the glass-substrate resonator mounting structure. The colour scale corresponds to the displacement from the unperturbed position of the structure (dashed contours). Mode labels correspond to those in \autoref{freqresponseFF}.}
	\label{FFmodes}
\end{figure}
The first observation is that the glass-substrate resonator mount exhibits fewer resonance peaks, especially within the \SIrange{0}{5}{\kilo\hertz} range, where we previously observed six modes. The lowest frequency mode for the glass-substrate resonator mount is located at $\SI{1.89}{\kilo\hertz}$ \textsc{(a)}, which is in sharp contrast to the fundamental bending mode at $\SI{267}{\hertz}$ for the fibre-tip resonator mount. Pushing the fundamental mode upward by a factor of seven reduces the structure's vulnerability to low-frequency ambient excitation. The proximity of the second mode at $\SI{1.98}{\kilo\hertz}$ to the first reflects the transition from the highly asymmetric rectangular cross-section to the more symmetric square cross-section of the current design, which is an unavoidable consequence of maximising the available steel volume within the (cuboid) glass cell. When comparing the effective modal mass ratios in \autoref{modeparticipationM} with those in \autoref{tab:modepercentFinalfinger}, we observe that the lowest two modes for the earlier design show a strong participation in the rotational degrees of freedom, with effective mass ratios of $\SI{85.1}{\percent}$ for $R_x$ (first mode) and $\SI{83.5}{\percent}$ for $R_z$ (second mode). In contrast, the glass-substrate resonator mount exhibits strong translational participation for the first two modes, with effective mass ratios of $\SI{35.8}{\percent}$ for $T_z$ (first mode) and $\SI{35.7}{\percent}$ for $T_x$ (second mode). This discrepancy arises because the fibre-tip resonator mount features a smaller support at the base, so the cantilever deforms by bending and rotating about an axis close to the origin. By contrast, the glass-substrate resonator mount consists of a tapered support that extends over nearly half the structure. As a result, the fundamental mode no longer pivots about an axis near the origin but instead buckles at the interface with the cuboid section of the mount, resulting in stronger modal mass participation for the translational degree of freedom.

To complement the previous analysis of structural acceleration, it is useful to consider the displacement response, since the cavity length is directly affected by deformations of the mount. This provides a more accurate representation of which modes are most detrimental to the stability of the resonator. Prominent acceleration peaks at high frequency may correspond to negligible displacements, in contrast to low-frequency modes. For a harmonic base acceleration $\vec{a}(t) = \vec{a}_0 \cos(\omega t)$, the corresponding rigid-body displacement amplitude scales as $\lvert \vec{x}(t) \rvert = \lvert \vec{a}_0 \rvert / \omega^2$. Thus, away from resonances, the displacement response decreases as $\sfrac{1}{\omega^2}$. The curves in \autoref{freqresponsecomparison} show the maximum displacement magnitude for both mounting designs under enforced base acceleration of $\lvert \vec{a}_0 \rvert = \SI{1}{\milli\metre\per\second\squared}$.
\begin{figure}[!t]
	%\vspace{0.7em}
	\centering
	\begin{tikzpicture}
	
	\begin{groupplot}[group style={group size=2 by 3, horizontal sep=0.5cm, vertical sep=0.5cm},xmin=0,ymin=0,height=4cm,width=7cm,no markers, ymin=1e-13, ymax=1e-5]
	
	\nextgroupplot
	[
	ymode=log,
	%ymin=1e-0,
	%ymax=1e4,
	%xlabel={Frequency},
	%ylabel={Max. linear acceleration},
	xmin=0, xmax=10,
	xticklabels={,,},
	%ymin=-0.1, ymax=1.1,
	%ymin=0, ymax=120,
	xtick={0,1,2,3,4,5,6,7,8,9,10},
	xticklabels={,,},
	ytick={1e-5,1e-6,1e-7,1e-8,1e-9,1e-10,1e-11,1e-12,1e-13},
	yticklabels={$10^{-5}$,,,,$10^{-9}$,,,,$10^{-13}$},
	%legend pos=north west,
	ymajorgrids=true,
	xmajorgrids=true,
	grid style=dashed,
	minor grid style=dashed,
	legend cell align={left},
	legend style={
		at={(1, 1)},  % Adjust (x, y) to move it
		anchor=north east,  % Attach the legend to its top-right corner
		draw=none,           % Optional: Remove border around the legend
		fill=none,
		align=left
	}
	]
	
	\addplot[
	color=markrood,
	name path=xpointone,
	] table [col sep = comma, x expr = \thisrowno{0}/1000, y expr = \thisrowno{1}/1000]{D_M_x_01.csv};
	
	\addplot[
	color=rainbow2of8,
	name path=xone,
	] table [col sep = comma, x expr = \thisrowno{0}/1000, y expr = \thisrowno{1}/1000]{D_M_x_1.csv};
	
	\addplot[
	color=rainbow4of8,
	name path=xfive,
	] table [col sep = comma, x expr = \thisrowno{0}/1000, y expr = \thisrowno{1}/1000]{D_M_x_5.csv};
	
	\addplot+[draw=none,domain=0:10,name path=xzero] {1e-15}; 
	
	\addplot[color=markrood, opacity=0.5] fill between[
	of = xone and xpointone,
	]; 
	
	\addplot[color=rainbow2of8, opacity=0.5] fill between[
	of = xfive and xone,
	]; 
	
	\addplot[color=rainbow4of8, opacity=0.5] fill between[
	of = xzero and xfive,
	]; 
	
	\pgfplotsset{
		after end axis/.append code={
			\node[align=center] at (rel axis cs: 0.5, 1.3) {\footnotesize\sc fibre-tip res.};
			\node[align=left, anchor=north west] at (rel axis cs: 0, 1) {\footnotesize $x$};
		}
	}
	
	\nextgroupplot
	[
	ymode=log,
	yticklabels={,,},
	%ymin=1e-0,
	%ymax=1e4,
	%xlabel={Frequency},
	%ylabel={Max. linear acceleration},
	xmin=0, xmax=10,
	xtick={0,1,2,3,4,5,6,7,8,9,10},
	xticklabels={,,},
	ytick={1e-5,1e-6,1e-7,1e-8,1e-9,1e-10,1e-11,1e-12,1e-13},
	yticklabels={,,},
	%ymin=-0.1, ymax=1.1,
	%ymin=0, ymax=120,
	%xtick={0,20,40,60,80,100},
	%ytick={0,20,40,60,80,100,120},
	%legend pos=north west,
	ymajorgrids=true,
	xmajorgrids=true,
	grid style=dashed,
	legend cell align={left},
	legend style={
		at={(1, 1)},  % Adjust (x, y) to move it
		anchor=north east,  % Attach the legend to its top-right corner
		draw=none,           % Optional: Remove border around the legend
		fill=none,
		align=left
	}
	]
	
	\addplot[
	color=markrood,
	name path=xpointone,
	] table [col sep = comma, x expr = \thisrowno{0}/1000, y expr = \thisrowno{1}/1000]{D_F_x_01.csv};
	
	\addplot[
	color=rainbow2of8,
	name path=xone,
	] table [col sep = comma, x expr = \thisrowno{0}/1000, y expr = \thisrowno{1}/1000]{D_F_x_1.csv};
	
	\addplot[
	color=rainbow4of8,
	name path=xfive,
	] table [col sep = comma, x expr = \thisrowno{0}/1000, y expr = \thisrowno{1}/1000]{D_F_x_5.csv};
	
	\addplot+[draw=none,domain=0:10,name path=xzero] {1e-15}; 
	
	\addplot[color=markrood, opacity=0.5] fill between[
	of = xone and xpointone,
	]; 
	
	\addplot[color=rainbow2of8, opacity=0.5] fill between[
	of = xfive and xone,
	]; 
	
	\addplot[color=rainbow4of8, opacity=0.5] fill between[
	of = xzero and xfive,
	]; 
	
	\addlegendentry{\footnotesize \hspace{0.3em} $G=\SI{0.1}{\percent}$}
	\addlegendentry{\footnotesize \hspace{0.3em} $G=\SI{1}{\percent}$}
	
	\pgfplotsset{
		after end axis/.append code={
			\node[align=center] at (rel axis cs: 0.5, 1.3) {\footnotesize\sc glass-substrate res.};
		}
	}
	
	\pgfplotsset{
		after end axis/.append code={
			\node[align=left, anchor=north west] at (rel axis cs: 0, 1) {\footnotesize $x$};
		}
	}
	
	\nextgroupplot
	[
	ymode=log,
	%ymin=1e-0,
	%ymax=1e4,
	%xlabel={Frequency},
	ylabel={Maximum displacement magnitude {[\si{\meter}]}},
	xmin=0, xmax=10,
	xtick={0,1,2,3,4,5,6,7,8,9,10},
	xticklabels={,,},
	ytick={1e-5,1e-6,1e-7,1e-8,1e-9,1e-10,1e-11,1e-12,1e-13},
	yticklabels={$10^{-5}$,,,,$10^{-9}$,,,,$10^{-13}$},
	%ymin=-0.1, ymax=1.1,
	%ymin=0, ymax=120,
	%xtick={0,20,40,60,80,100},
	%ytick={0,20,40,60,80,100,120},
	%legend pos=north west,
	ymajorgrids=true,
	xmajorgrids=true,
	grid style=dashed,
	legend cell align={left},
	legend style={
		at={(1, 1)},  % Adjust (x, y) to move it
		anchor=north east,  % Attach the legend to its top-right corner
		draw=none,           % Optional: Remove border around the legend
		fill=none,
		align=left
	}
	]
	
	\addplot[
	color=markrood,
	name path=xpointone,
	] table [col sep = comma, x expr = \thisrowno{0}/1000, y expr = \thisrowno{1}/1000]{D_M_y_01.csv};
	
	\addplot[
	color=rainbow2of8,
	name path=xone,
	] table [col sep = comma, x expr = \thisrowno{0}/1000, y expr = \thisrowno{1}/1000]{D_M_y_1.csv};
	
	\addplot[
	color=rainbow4of8,
	name path=xfive,
	] table [col sep = comma, x expr = \thisrowno{0}/1000, y expr = \thisrowno{1}/1000]{D_M_y_5.csv};
	
	\addplot+[draw=none,domain=0:10,name path=xzero] {1e-15}; 
	
	\addplot[color=markrood, opacity=0.5] fill between[
	of = xone and xpointone,
	]; 
	
	\addplot[color=rainbow2of8, opacity=0.5] fill between[
	of = xfive and xone,
	]; 
	
	\addplot[color=rainbow4of8, opacity=0.5] fill between[
	of = xzero and xfive,
	]; 
	
	\pgfplotsset{
		after end axis/.append code={
			\node[align=left, anchor=north west] at (rel axis cs: 0, 1) {\footnotesize $y$};
		}
	}
	
	\nextgroupplot
	[
	ymode=log,
	yticklabels={,,},
	%ymin=1e-0,
	%ymax=1e4,
	%xlabel={Frequency},
	%ylabel={Max. linear acceleration},
	xmin=0, xmax=10,
	xtick={0,1,2,3,4,5,6,7,8,9,10},
	xticklabels={,,},
	ytick={1e-5,1e-6,1e-7,1e-8,1e-9,1e-10,1e-11,1e-12,1e-13},
	yticklabels={,,},
	%ymin=-0.1, ymax=1.1,
	%ymin=0, ymax=120,
	%xtick={0,20,40,60,80,100},
	%ytick={0,20,40,60,80,100,120},
	%legend pos=north west,
	ymajorgrids=true,
	xmajorgrids=true,
	grid style=dashed,
	legend cell align={left},
	legend style={
		at={(1, 1)},  % Adjust (x, y) to move it
		anchor=north east,  % Attach the legend to its top-right corner
		draw=none,           % Optional: Remove border around the legend
		fill=none,
		align=left
	}
	]
	
	\addplot[
	color=markrood,
	name path=xpointone,
	] table [col sep = comma, x expr = \thisrowno{0}/1000, y expr = \thisrowno{1}/1000]{D_F_y_01.csv};
	
	\addplot[
	color=rainbow2of8,
	name path=xone,
	] table [col sep = comma, x expr = \thisrowno{0}/1000, y expr = \thisrowno{1}/1000]{D_F_y_1.csv};
	
	\addplot[
	color=rainbow4of8,
	name path=xfive,
	] table [col sep = comma, x expr = \thisrowno{0}/1000, y expr = \thisrowno{1}/1000]{D_F_y_5.csv};
	
	\addplot+[draw=none,domain=0:10,name path=xzero] {1e-15}; 
	
	\addplot[color=markrood, opacity=0.5] fill between[
	of = xone and xpointone,
	]; 
	
	\addplot[color=rainbow2of8, opacity=0.5] fill between[
	of = xfive and xone,
	]; 
	
	\addplot[color=rainbow4of8, opacity=0.5] fill between[
	of = xzero and xfive,
	]; 
	
	\pgfplotsset{
		after end axis/.append code={
			\node[align=left, anchor=north west] at (rel axis cs: 0, 1) {\footnotesize $y$};
		}
	}
	
	\addlegendentry{\footnotesize \hspace{0.3em} $G=\SI{5}{\percent}$}
	
	\nextgroupplot
	[
	ymode=log,
	%ymin=1e-0,
	%ymax=1e4,
	%xlabel={Frequency},
	%ylabel={Max. linear acceleration},
	xmin=0, xmax=10,
	xtick={0,1,2,3,4,5,6,7,8,9,10},
	xticklabels={$0$,,$2$,,$4$,,$6$,,$8$,,$10$},
	ytick={1e-5,1e-6,1e-7,1e-8,1e-9,1e-10,1e-11,1e-12,1e-13},
	yticklabels={$10^{-5}$,,,,$10^{-9}$,,,,$10^{-13}$},
	%xticklabels={,,},
	%ymin=-0.1, ymax=1.1,
	%ymin=0, ymax=120,
	%xtick={0,20,40,60,80,100},
	%ytick={0,20,40,60,80,100,120},
	%legend pos=north west,
	ymajorgrids=true,
	xmajorgrids=true,
	grid style=dashed,
	legend cell align={left},
	legend style={
		at={(1, 1)},  % Adjust (x, y) to move it
		anchor=north east,  % Attach the legend to its top-right corner
		draw=none,           % Optional: Remove border around the legend
		fill=none,
		align=left
	}
	]
	
	\addplot[
	color=markrood,
	name path=xpointone,
	] table [col sep = comma, x expr = \thisrowno{0}/1000, y expr = \thisrowno{1}/1000]{D_M_z_01.csv};
	
	\addplot[
	color=rainbow2of8,
	name path=xone,
	] table [col sep = comma, x expr = \thisrowno{0}/1000, y expr = \thisrowno{1}/1000]{D_M_z_1.csv};
	
	\addplot[
	color=rainbow4of8,
	name path=xfive,
	] table [col sep = comma, x expr = \thisrowno{0}/1000, y expr = \thisrowno{1}/1000]{D_M_z_5.csv};
	
	\addplot+[draw=none,domain=0:10,name path=xzero] {1e-15}; 
	
	\addplot[color=markrood, opacity=0.5] fill between[
	of = xone and xpointone,
	]; 
	
	\addplot[color=rainbow2of8, opacity=0.5] fill between[
	of = xfive and xone,
	]; 
	
	\addplot[color=rainbow4of8, opacity=0.5] fill between[
	of = xzero and xfive,
	]; 
	
	\pgfplotsset{
		after end axis/.append code={
			\node[align=left, anchor=north west] at (rel axis cs: 0, 1) {\footnotesize $z$};
		}
	}
	
	\nextgroupplot
	[
	ymode=log,
	xtick={0,1,2,3,4,5,6,7,8,9,10},
	xticklabels={$0$,,$2$,,$4$,,$6$,,$8$,,$10$},
	ytick={1e-5,1e-6,1e-7,1e-8,1e-9,1e-10,1e-11,1e-12,1e-13},
	yticklabels={,,},
	%ymin=1e-0,
	%ymax=1e4,
	xlabel={Enforced acceleration frequency {[\si{\kilo\hertz}]}},
	xlabel style={shift={(-7em, 0em)}},
	%ylabel={Max. linear acceleration},
	xmin=0, xmax=10,
	%xticklabels={,,},
	%ymin=-0.1, ymax=1.1,
	%ymin=0, ymax=120,
	%xtick={0,20,40,60,80,100},
	%ytick={0,20,40,60,80,100,120},
	%legend pos=north west,
	ymajorgrids=true,
	xmajorgrids=true,
	grid style=dashed,
	legend cell align={left},
	legend style={
		at={(1, 1)},  % Adjust (x, y) to move it
		anchor=north east,  % Attach the legend to its top-right corner
		draw=none,           % Optional: Remove border around the legend
		fill=none,
		align=left
	}
	]
	
	\addplot[
	color=markrood,
	name path=xpointone,
	] table [col sep = comma, x expr = \thisrowno{0}/1000, y expr = \thisrowno{1}/1000]{D_F_z_01.csv};
	
	\addplot[
	color=rainbow2of8,
	name path=xone,
	] table [col sep = comma, x expr = \thisrowno{0}/1000, y expr = \thisrowno{1}/1000]{D_F_z_1.csv};
	
	\addplot[
	color=rainbow4of8,
	name path=xfive,
	] table [col sep = comma, x expr = \thisrowno{0}/1000, y expr = \thisrowno{1}/1000]{D_F_z_5.csv};
	
	\addplot+[draw=none,domain=0:10,name path=xzero] {1e-15}; 
	
	\addplot[color=markrood, opacity=0.5] fill between[
	of = xone and xpointone,
	]; 
	
	\addplot[color=rainbow2of8, opacity=0.5] fill between[
	of = xfive and xone,
	]; 
	
	\addplot[color=rainbow4of8, opacity=0.5] fill between[
	of = xzero and xfive,
	]; 
	
	\pgfplotsset{
		after end axis/.append code={
			\node[align=left, anchor=north west] at (rel axis cs: 0, 1) {\footnotesize $z$};
		}
	}
	
	
	\end{groupplot}
	
	
	
	\end{tikzpicture}
	\caption[Comparison of maximum displacement from the frequency response analyses of two resonator mount designs]{Comparison of maximum displacement from the frequency response analyses of two resonator mount designs. Acceleration is enforced at the base of each mount is in three different directions: $x$ \textsc{(top)}, $y$ \textsc{(middle)} and $z$ \textsc{(bottom)}. In each direction, traces are shown for three different values of the structural damping coefficient $G$. The magnitude of the enforced acceleration equals \SI{1}{\milli\meter\per\second\squared}}
	\label{freqresponsecomparison} 
\end{figure}
The comparison highlights clear advantages of the glass-substrate resonator mount. For enforced acceleration along the $x$-axis, the first visible displacement resonance is shifted upward by about \SI{1}{\kilo\hertz} and reduced in magnitude relative to the fibre-tip resonator mount, with all higher modes remaining below \SI{1}{\nano\metre} of maximum displacement. Along the $y$-axis, the response at resonance is suppressed by roughly an order of magnitude across the spectrum. Most strikingly, for excitation along the $z$-axis the problematic \SI{267}{\hertz} peak of the fibre-tip resonator mount is eliminated. The most prominent resonance in this direction now occurs at \SI{1.89}{\kilo\hertz}, but as the first visible mode its peak value is reduced by nearly two orders of magnitude. All other resonances again lie below \SI{1}{\nano\metre} of maximum displacement. For the current design, this means that only the closely spaced pair of modes at \SI{1.89}{\kilo\hertz} and \SI{1.98}{\kilo\hertz} remain of concern. Their isolation, however, may even permit a split feedback strategy in which one piezoelectric actuator suppresses sub-\SI{1.89}{\kilo\hertz} disturbances and another covers the band above \SI{1.98}{\kilo\hertz}. Since feedback gain must be reduced in the vicinity of sharp resonances to prevent instability, separating the control bands in this way avoids having either actuator operate across these problematic frequencies. With broader resonance-free regions and substantially reduced displacement amplitudes, the current mounting design provides a more reliable platform for the stable operation of an optical resonator.

\section{Pyramidal magneto-optical trap}
\label{sec:pyramidal_mot}
asd asdfa sdfa sdf asf asf asdf asdf asdf asdf asdf asdf dsasd asdfa sdfa sdf asf asf asdf asdf asdf asdf asdf asdf dsasd asdfa sdfa sdf asf asf asdf asdf asdf asdf asdf asdf dsasd asdfa sdfa sdf asf asf asdf asdf asdf asdf asdf asdf dsasd asdfa sdfa sdf asf asf asdf asdf asdf asdf asdf asdf dsasd asdfa sdfa sdf asf asf asdf asdf asdf asdf asdf asdf dsasd asdfa sdfa sdf asf asf asdf asdf asdf asdf asdf asdf dsasd asdfa sdfa sdf asf asf asdf asdf asdf asdf asdf asdf ds

\iffalse
By securing these elements on a single mount, the architecture minimises mechanical degrees of freedom, reduces alignment complexity and ensures stable spatial overlap between the atomic source, the dipole trap and the cavity mode.


This arrangement reflects a single incident cooling beam into the standard six-beam MOT configuration, providing efficient loading of the dipole trap from a cold atomic ensemble. 




The choice of a pyramidal MOT as the cold-atom source is motivated by its compactness, mechanical robustness, and minimal optical access requirements. Conventional MOTs use three orthogonal pairs of counterpropagating laser beams, requiring multiple independent optical elements and substantial optical access to the trapping region. While retroreflection can halve the number of independent beams, the apparatus remains relatively bulky and alignment-sensitive. In contrast, the pyramidal MOT uses a set of four mirrors in a square-based pyramidal configuration to reflect a single circularly polarised input beam into the full six-beam arrangement required for magneto-optical trapping in a spherical quadrupole magnetic field. Each pair of counterpropagating beams emerges with the correct relative circular polarisation by virtue of the $180^{\circ}$ phase shift between $s$- and $p$-components upon reflection from the mirrors. 

This geometry yields balanced beam intensities at all points in the trapping volume (assuming well-aligned and identical mirrors), making efficient use of the available optical power and suppressing phase and intensity noise via common-mode rejection. Moreover, because the light returns along the incident axis, only a single optical feedthrough is required, allowing the MOT optics to be integrated entirely within the vacuum chamber. By truncating the pyramid apex to form an aperture, the configuration can be operated as a \emph{source}, in which atoms are extracted along the symmetry axis to load a downstream trap---in our case, the optical dipole trap inside the cavity mode.

Integrating the pyramidal MOT directly into the cavity mounting finger offers two primary benefits: (i) the MOT is mechanically referenced to the cavity, ensuring that any long-term drifts affect both subsystems identically; and (ii) the use of a single, compact mechanical support minimises vibrational resonances and simplifies the optical layout. The remainder of this chapter describes the design considerations, mechanical analysis, and optical performance of the integrated cavity support and MOT source, beginning with the cavity mounting finger design and followed by the realisation and characterisation of the pyramidal MOT.

\fi

\subsection{Overview}

5 - 2 * ((6. * a * t) ** (1./3.))


\begin{figure}[t]
	\hspace{-0cm}
	\begin{tikzpicture}[remember picture]
		\node[anchor=south west,inner sep=0] (image1) at (0,0) {\includegraphics[width=1\textwidth]{{"pyramidinfinger"}.pdf}};
		
		% left four
		\begin{scope}[x={(image1.south east)},y={(image1.north west)}]
			% draw a grid
			%\draw[help lines,xstep=.1,ystep=.1,overlay] (0,0) grid (1,1);
			% draw ticks
			%\foreach \x in {0,1,...,9} { \node [anchor=north,overlay] at (\x/10,0) {0.\x}; }
			%\foreach \y in {0,1,...,9} { \node [anchor=east,overlay] at (0,\y/10) {0.\y}; }
			
			% labels
			%\node[markgrijs, align=center, rotate=0] at (0.0,0.9) {\footnotesize\sc (a)};
			%\node[markgrijs, align=center, rotate=0] at (0.35,0.9) {\footnotesize\sc (b)};
			%\node[markgrijs, align=center, rotate=0] at (0.7,0.9) {\footnotesize\sc (c)};
			
			\node[align=center, rotate=0] at (0.15,-0.05) {\footnotesize\sc (a)};
			\node[align=center, rotate=0] at (0.5,-0.05) {\footnotesize\sc (b)};
			\node[align=center, rotate=0] at (0.85,-0.05) {\footnotesize\sc (c)};
			
		\end{scope}
		
	\end{tikzpicture}
	\caption[Loading an intracavity dipole trap from a pyramidal MOT]{Loading an intracavity dipole trap from a pyramidal MOT. \textsc{(a)} A cutaway view of the cavity mounting structure, revealing the cavity mirrors and part of the right-angle mirrors that reflect the cooling beams to form the pyramidal MOT. At the apex of the pyramid, a central opening allows an atomic beam to emanate towards the cavity mode. \textsc{(b)} The same structure viewed from the back, showing the gold-coated rear surface of the mirror holder. The location of the cross-section shown in \textsc{(c)} is indicated in light blue. \textsc{(c)} A cross-sectional view illustrating the dipole trap beam (red) and the cooling beams (green), which are incident at small angles relative to the normal. These beams reflect off the gold coating, creating a region around the dipole trap (yellow) where an optical molasses can form, populated by cold atoms from the pyramidal MOT.}
	\label{pyramidloading} 
\end{figure}

\subsection{Aperture diffraction analysis}

\begin{figure}[!t]
	\centering
	\begin{tikzpicture}[remember picture]
	\node[anchor=south west,inner sep=0] (image1) at (0,0) {\includegraphics[width=0.75\textwidth]{{"motdiffraction"}.pdf}};
	\begin{scope}[x={(image1.south east)},y={(image1.north west)}]
	% draw a grid
	%\draw[help lines,xstep=.1,ystep=.1,overlay] (0,0) grid (1,1);
	% draw ticks
	%\foreach \x in {0,1,...,9} { \node [anchor=north,overlay] at (\x/10,0) {0.\x}; }
	%\foreach \y in {0,1,...,9} { \node [anchor=east,overlay] at (0,\y/10) {0.\y}; }
	
	\node[markgrijs, align=center, rotate=0] at (0.13,-0.1) {\footnotesize\sc front view}; 
	\node[markgrijs, align=center, rotate=0] at (0.45,-0.1) {\footnotesize\sc side view}; 
	
	\node[white, align=center, rotate=0, anchor=center] at (0.52,0.5) {\footnotesize \SI{780}{\nano\meter}}; 
	\node[white, align=center, rotate=0, anchor=center] at (0.73,0.5) {\footnotesize \SI{780}{\nano\meter}}; 
	
	\draw[markgrijs, ->] (0.848, 0.2) -- (0.848, 0.4) node[left, pos=0.5] {\footnotesize $a$};
	\draw[markgrijs, ->] (0.848, 0.8) -- (0.848, 0.6) node[above, anchor=mid] {\footnotesize };
	
	\draw[markgrijs, dotted] (0.848, 0.2) -- (0.848, 0.05) node[above, anchor=mid] {};
	\draw[markgrijs, <->] (0.848, 0.1) -- (0.983, 0.1) node[below, pos=0.5] {\footnotesize $d$};
	
	\node[markgrijs, align=center, anchor=center] at (0.6245, 0.5) {\footnotesize $=$};
	
	\end{scope}
	
	\end{tikzpicture}
	\caption[Schematic of the pyramidal MOT mirror configuration with an apex aperture and corresponding diffraction model]{Schematic of the pyramidal MOT mirror configuration with an aperture at the apex. The front view shows the four reflective faces around the square central aperture. The side view illustrates the incident and reflected \SI{780}{\nano\meter} laser beam. The diffraction of the reflected beam due to the finite aperture size is equivalently described as diffraction from a square obstruction of size $a$ (right). The diffraction pattern is modelled at various distances $d$.}
	\label{motdiffraction} 
\end{figure}

\begin{figure}[t]
	\centering
	\begin{tikzpicture}
	
	\begin{groupplot}[group style={group size=8 by 8,horizontal sep=0.1cm, vertical sep=0.1cm},height=3cm,width=3cm,no markers]
	
	\pgfplotsinvokeforeach{1,...,64} {
		\nextgroupplot[
		ylabel={},
		xmin=0, xmax=1,
		ymin=0, ymax=1,
		xticklabels={,\textcolor{white}{0},},
		yticklabels={,,},
		ymajorgrids=true,
		xmajorgrids=true,
		grid style=dashed,
		axis lines=none,
		xlabel={},
		ylabel={},
		]
		\addplot graphics [
		xmin=0, xmax=1,
		ymin=0, ymax=1
		] {motdiff#1.pdf}; % #1 is the current number
		
		\coordinate (topnode#1) at (axis cs: -0.1,1.1);
		\coordinate (bottomnode#1) at (axis cs: -0.1,-0.15);
		\coordinate (rightnode#1) at (axis cs: 1.15,1.1);
		
		\coordinate (vert#1) at (axis cs: -0.35,0.5);
		\coordinate (hor#1) at (axis cs: 0.5,1.25);
		
		\coordinate (edgeleft#1) at (axis cs: 0.11,0.5);
		\coordinate (edgeright#1) at (axis cs: 0.89,0.5);
		
		\coordinate (fresnel#1) at (axis cs: 0.5,0.9);
		\coordinate (mfresnel#1) at (axis cs: 0.5,0.5);
	}
	

	
	\end{groupplot}
	
	% Set constants
	\pgfmathsetmacro{\lambda}{0.780} % Wavelength in millimetres
	
	% Fresnel numbers for combinations of a and d
	% The pattern: a from {0.1,...,0.8}, d from {0,...,7} → 64 values total
	\pgfplotsinvokeforeach{1,...,64}{
		
		% Convert #1 to indices for a and d:
		% a-index: floor((#1-1)/8), d-index: mod(#1-1,8)
		\pgfmathtruncatemacro{\aindex}{int((#1-1)/8)}
		\pgfmathtruncatemacro{\dindex}{mod(#1-1,8)}
		
		% Map indices to actual values:
		\pgfmathsetmacro{\a}{0.1*(\aindex+1)} % a in {0.1,...,0.8}
		\pgfmathsetmacro{\d}{\dindex}          % d in {0,...,7}
		
		% Calculate Fresnel number:
		\ifnum\dindex=0
			\pgfmathsetmacro{\result}{0}
		\else
			\pgfmathsetmacro{\boolean}{ifthenelse(\a<0.5, 0, 1)} 
			\ifnum\boolean=0
				\pgfmathsetmacro{\result}{1e3*(\a^2)/(\lambda*\d)}
				\node[align=center, white] at (fresnel#1) {\tiny $\pgfmathprintnumber[precision=1]{\result}$};
			\else
				\pgfmathsetmacro{\result}{1e3*(\a^2)/(\lambda*\d)}
				\node[align=center, white] at (mfresnel#1) {\tiny $\pgfmathprintnumber[precision=1]{\result}$};
			\fi
		\fi
	}
	
	
	% Print with three significant digits
	
	

	
	\draw[white, <->] (edgeleft57) -- (edgeright57) node[below, pos=0.5] {\small $a$};
	
	\draw[->, line cap=rect] (topnode1) -- (bottomnode57) node[left, rotate=0] {$a$};
	\draw[->, line cap=rect] (topnode1) -- (rightnode8) node[above] {\small $d$};
	
	
	\node[align=center] at (vert1) {\small $0.1$};
	\node[align=center] at (vert9) {\small $0.2$};
	\node[align=center] at (vert17) {\small $0.3$};
	\node[align=center] at (vert25) {\small $0.4$};
	\node[align=center] at (vert33) {\small $0.5$};
	\node[align=center] at (vert41) {\small $0.6$};
	\node[align=center] at (vert49) {\small $0.7$};
	\node[align=center] at (vert57) {\small $0.8$};
	
	\node[align=center] at (hor1) {\small $0$};
	\node[align=center] at (hor2) {\small $1$};
	\node[align=center] at (hor3) {\small $2$};
	\node[align=center] at (hor4) {\small $3$};
	\node[align=center] at (hor5) {\small $4$};
	\node[align=center] at (hor6) {\small $5$};
	\node[align=center] at (hor7) {\small $6$};
	\node[align=center] at (hor8) {\small $7$};
	
	\end{tikzpicture}
	\caption[Simulated diffraction patterns from a pyramid MOT mirror configuration]{Simulated diffraction patterns from a pyramid MOT mirror configuration with a square aperture of various dimensions $a$ at the apex, observed at distances $d$. The Fresnel number $F$ is indicated in each subplot. Each pattern represents a $\SI{1}{\milli\meter}\times\SI{1}{\milli\meter}$ area in the observation plane.}
\label{diffraction2d}
\end{figure}




\begin{figure}[t]
	\centering
	\begin{tikzpicture}
	
	\pgfplotstableread[col sep=comma]{differrs.csv}\datatable
	\pgfplotstableread[col sep=comma]{diffmaxerrs2d.csv}\datatablemax
		
		\begin{groupplot}[group style={group size=8 by 8,horizontal sep=0.1cm, vertical sep=0.1cm},height=3cm,width=3cm,no markers]
		
		\pgfplotsinvokeforeach{1,...,64} {
			%\pgfmathtruncatemacro{\rowindex}{1+floor((#1-1)/8)}
			% Plots with indices 8,16,24,...,64 -> show y-tick labels on the right
			\nextgroupplot[
			ylabel={},
			ymin=0, ymax=2,
			xmin=0, xmax=1,
			xticklabels={,\textcolor{white}{0},},
			yticklabels={,,},
			yticklabel pos=right,
			y axis line style={draw=black},
			ymajorgrids=false,
			xmajorgrids=false,
			grid style=dashed,
			xlabel={},
			xtick=\empty,
			ytick=\empty,
			]
			
			
			% Load data from CSV files (motdiffcs#.csv with x,y columns)
			\addplot[
			name path=curve,
			color=markrood,
			] table [x index=0, y index=1, col sep=comma] {motdiffcs#1.csv};
			
			\addplot[
			name path=axis,
			draw=none
			] coordinates {(0,0) (1,0)};
			
			\addplot[
			fill=markrood,
			fill opacity=0.3
			] fill between [of=curve and axis];
			
			\draw[dashed, markgrijs] (0,1) -- (1,1);
			
			\coordinate (topnode#1) at (axis cs:-0.1,2.2);
			\coordinate (bottomnode#1) at (axis cs:-0.1,-0.3);
			\coordinate (rightnode#1) at (axis cs:1.15,2.2);
			
			\coordinate (vert#1) at (axis cs:-0.35,1);
			\coordinate (hor#1) at (axis cs:0.5,2.5);
			
			\coordinate (edgeleft#1) at (axis cs:0.11,0.7);
			\coordinate (edgeright#1) at (axis cs:0.89,0.7);
			
			\coordinate (hedgeleft#1) at (axis cs:0,1.7);
			\coordinate (hedgeright#1) at (axis cs:1,1.7);
			
			\coordinate (outerright#1) at (axis cs:1.15,1);
			\coordinate (louterright#1) at (axis cs:1.15,0);
			
			\pgfmathtruncatemacro{\aindex}{int((#1-1)/8)}
			\pgfmathtruncatemacro{\dindex}{mod(#1-1,8)}
			
			% Map indices to actual values:
			\pgfmathsetmacro{\a}{0.1*(\aindex+1)} % a in {0.1,...,0.8}
			\pgfmathsetmacro{\d}{\dindex}          % d in {0,...,7}
			
			% Calculate Fresnel number:
			\ifnum\dindex=0
				\pgfmathsetmacro{\result}{0}
			\else
				\pgfmathparse{#1 - 1}
				\pgfplotstablegetelem{\pgfmathresult}{[index] 0}\of\datatable % Get nth row, 1st column
				\edef\valueone{\pgfplotsretval} % Store value
				\pgfmathsetmacro{\resultone}{\valueone}
				\node[align=left, markgrijs, anchor=west] at (axis cs: -0.05, 1.8) {\tiny $\pgfmathprintnumber[precision=1]{\resultone}$};
				
				\pgfplotstablegetelem{\pgfmathresult}{[index] 0}\of\datatablemax % Get nth row, 1st column
				\edef\valuetwo{\pgfplotsretval} % Store value
				\pgfmathsetmacro{\resulttwo}{\valuetwo*100}
				\node[align=right, markgrijs, anchor=east] at (axis cs: 1.05, 1.8) {\tiny $\pgfmathprintnumber[precision=0]{\resulttwo}$};
			\fi
		}
		
		\end{groupplot}
		
		
		
		\draw[<->] (edgeleft57) -- (edgeright57) node[below, pos=0.5] {\footnotesize $a$};
		\draw[<->] (hedgeleft57) -- (hedgeright57) node[below, pos=0.5] {\footnotesize \SI{1}{\milli\meter}};
		
		\draw[->, line cap=rect] (topnode1) -- (bottomnode57) node[left, rotate=0] {\small $a$};
		\draw[->, line cap=rect] (topnode1) -- (rightnode8) node[above] {\small $d$};
		
		
		\node[align=center] at (vert1) {\small $0.1$};
		\node[align=center] at (vert9) {\small $0.2$};
		\node[align=center] at (vert17) {\small $0.3$};
		\node[align=center] at (vert25) {\small $0.4$};
		\node[align=center] at (vert33) {\small $0.5$};
		\node[align=center] at (vert41) {\small $0.6$};
		\node[align=center] at (vert49) {\small $0.7$};
		\node[align=center] at (vert57) {\small $0.8$};
		
		\node[align=center] at (hor1) {\small $0$};
		\node[align=center] at (hor2) {\small $1$};
		\node[align=center] at (hor3) {\small $2$};
		\node[align=center] at (hor4) {\small $3$};
		\node[align=center] at (hor5) {\small $4$};
		\node[align=center] at (hor6) {\small $5$};
		\node[align=center] at (hor7) {\small $6$};
		\node[align=center] at (hor8) {\small $7$};
		
		\node[markgrijs, align=center] at (outerright8) {\small $1$};
		\node[markgrijs, align=center] at (outerright16) {\small $1$};
		\node[markgrijs, align=center] at (outerright24) {\small $1$};
		\node[markgrijs, align=center] at (outerright32) {\small $1$};
		\node[markgrijs, align=center] at (outerright40) {\small $1$};
		\node[markgrijs, align=center] at (outerright48) {\small $1$};
		\node[markgrijs, align=center] at (outerright56) {\small $1$};
		\node[markgrijs, align=center] at (outerright64) {\small $1$};
		
		\node[markgrijs, align=center] at (louterright8) {\small $0$};
		\node[markgrijs, align=center] at (louterright16) {\small $0$};
		\node[markgrijs, align=center] at (louterright24) {\small $0$};
		\node[markgrijs, align=center] at (louterright32) {\small $0$};
		\node[markgrijs, align=center] at (louterright40) {\small $0$};
		\node[markgrijs, align=center] at (louterright48) {\small $0$};
		\node[markgrijs, align=center] at (louterright56) {\small $0$};
		\node[markgrijs, align=center] at (louterright64) {\small $0$};
		
	\end{tikzpicture}
	\caption[Central cross-sections of the diffraction profiles from a pyramid MOT mirror configuration]{Central cross-sections of the diffraction profiles from \autoref{diffraction2d}, showing the normalised intensity (such that $1$ corresponds to the maximum incident intensity) along a \SI{1}{\milli\meter} region in front of the aperture for various values of $a$ and $d$. Within each plot, the top left number (with arbitrary units) represents the integrated absolute difference from the corresponding cross-section at $d=0$ (integrated over a distance of \SI{10}{\milli\meter}). The top right number is the maximum absolute difference as a percentage of maximum incident intensity.}
\label{diffraction1d}
\end{figure}




\begin{figure}[t]
	\hspace{-0.15cm}
	\begin{tikzpicture}
	\begin{groupplot}[group style={group size=3 by 4,horizontal sep=0.9cm, vertical sep=1.7cm},height=5cm,width=5cm,no markers]
	
		\nextgroupplot
	[
	%xlabel={$x$},
	%ylabel={$y$},
	enlargelimits=false,
	%xmin=0.5,
	%xmax=4.5,
	%ymin=0.5,
	%ymax=4.5,
	axis on top,    % <-- (use with `matrix plot*')
	view={0}{90},
	colormap/viridis,
	colorbar horizontal,
	colorbar style={
		at={(0,1.2)},
		anchor=north west,
		xticklabel pos=upper,
		%x tick scale label style={yshift=-25pt, xshift=23pt},
	},
	point meta=explicit,
	point meta min=0.0,
	point meta max=2.0,
	%xtick={1,2,3,4},
	%ytick={1,2,3,4},
	%xticklabels={1,2,3,4},
	%yticklabels={1,2,3,4},
	]
	\addplot graphics [
	xmin=0, xmax=10,  % Stretch to fit the x-axis limits
	ymin=0, ymax=10    % Stretch to fit the y-axis limits
	] {fresnelfull.pdf};
	
	\draw[fill=none, white] (axis cs:4.5,4.5) rectangle (axis cs:5.5,5.5);
	
	\pgfplotsset{
	after end axis/.append code={
		\node[align=center] at (axis cs: 5, 15) {\footnotesize\sc Fresnel approx.};
		}
	}
	
	
	
	\nextgroupplot
	[
	%xlabel={$x$},
	%ylabel={$y$},
	enlargelimits=false,
	%xmin=0.5,
	%xmax=4.5,
	%ymin=0.5,
	%ymax=4.5,
	axis on top,    % <-- (use with `matrix plot*')
	view={0}{90},
	colormap/viridis,
	colorbar horizontal,
	colorbar style={
		at={(0,1.2)},
		anchor=north west,
		xticklabel pos=upper,
		%x tick scale label style={yshift=-25pt, xshift=23pt},
	},
	point meta=explicit,
	point meta min=0.0,
	point meta max=2.0,
	%xtick={1,2,3,4},
	%ytick={1,2,3,4},
	%xticklabels={1,2,3,4},
	%yticklabels={1,2,3,4},
	]
	\addplot graphics [
	xmin=0, xmax=10,  % Stretch to fit the x-axis limits
	ymin=0, ymax=10    % Stretch to fit the y-axis limits
	] {kirfull.pdf};
	
	\draw[fill=none, white] (axis cs:4.5,4.5) rectangle (axis cs:5.5,5.5);
	
	\pgfplotsset{
	after end axis/.append code={
		\node[align=center] at (axis cs: 5, 15) {\footnotesize\sc direct integration};
		}
	}
	
	
		\nextgroupplot
	[
	%xlabel={$x$},
	%ylabel={$y$},
	enlargelimits=false,
	%xmin=0.5,
	%xmax=4.5,
	%ymin=0.5,
	%ymax=4.5,
	axis on top,    % <-- (use with `matrix plot*')
	view={0}{90},
	colormap/viridis,
	colorbar horizontal,
	colorbar style={
		at={(0,1.2)},
		anchor=north west,
		xticklabel pos=upper,
		xtick={0,2.5,5,7.5,10},
		xticklabels={0,2.5,5,7.5,10}
		%x tick scale label style={yshift=-25pt, xshift=23pt},
	},
	point meta=explicit,
	point meta min=0.0,
	point meta max=10,
	%xtick={1,2,3,4},
	%ytick={1,2,3,4},
	%xticklabels={1,2,3,4},
	%yticklabels={1,2,3,4},
	]
	\addplot graphics [
	xmin=0, xmax=10,  % Stretch to fit the x-axis limits
	ymin=0, ymax=10    % Stretch to fit the y-axis limits
	] {errorfull.pdf};
	
	\draw[fill=none, white] (axis cs:4.5,4.5) rectangle (axis cs:5.5,5.5);
	
	\pgfplotsset{
		after end axis/.append code={
			\node[align=center] at (axis cs: 5, 15) {\footnotesize\sc error {[\si{\percent}]}};
		}
	}
	
	
	
	\nextgroupplot
	[
	%xlabel={$x$},
	ylabel={$y$ {[\si{\milli\meter}]}},
	ylabel style={xshift=2.5cm},
	enlargelimits=false,
	%xmin=0.5,
	%xmax=4.5,
	%ymin=0.5,
	%ymax=4.5,
	axis on top,    % <-- (use with `matrix plot*')
	view={0}{90},
	colormap/viridis,
	colorbar horizontal,
	colorbar style={
		at={(0,1.2)},
		anchor=north west,
		xticklabel pos=upper,
		%x tick scale label style={yshift=-25pt, xshift=23pt},
	},
	point meta=explicit,
	point meta min=0.0,
	point meta max=2.0,
	%xtick={1,2,3,4},
	%ytick={1,2,3,4},
	%xticklabels={1,2,3,4},
	%yticklabels={1,2,3,4},
	]
	\addplot graphics [
	xmin=0, xmax=1,  % Stretch to fit the x-axis limits
	ymin=0, ymax=1    % Stretch to fit the y-axis limits
	] {fresnelzoom.pdf};
	
	\addplot graphics [
	xmin=0.6, xmax=0.98,  % Stretch to fit the x-axis limits
	ymin=0.6, ymax=0.98    % Stretch to fit the y-axis limits
	] {fresnelsmall.pdf};
	
	\draw[fill=none, white] (axis cs:0.6,0.6) rectangle (axis cs:0.98,0.98);
	
	
	\node (lowerleft) at (axis cs: 0, -1.5) {};
	
	
		\nextgroupplot
	[
	xlabel={$x$ {[\si{\milli\meter}]}},
	%ylabel={$y$},
	enlargelimits=false,
	%xmin=0.5,
	%xmax=4.5,
	%ymin=0.5,
	%ymax=4.5,
	axis on top,    % <-- (use with `matrix plot*')
	view={0}{90},
	colormap/viridis,
	colorbar horizontal,
	colorbar style={
		at={(0,1.2)},
		anchor=north west,
		xticklabel pos=upper,
		%x tick scale label style={yshift=-25pt, xshift=23pt},
	},
	point meta=explicit,
	point meta min=0.0,
	point meta max=2.0,
	%xtick={1,2,3,4},
	%ytick={1,2,3,4},
	%xticklabels={1,2,3,4},
	%yticklabels={1,2,3,4},
	]
	\addplot graphics [
	xmin=0, xmax=1,  % Stretch to fit the x-axis limits
	ymin=0, ymax=1    % Stretch to fit the y-axis limits
	] {kirzoom.pdf};

	\addplot graphics [
	xmin=0.6, xmax=0.98,  % Stretch to fit the x-axis limits
	ymin=0.6, ymax=0.98    % Stretch to fit the y-axis limits
	] {kirsmall.pdf};
	
	\draw[fill=none, white] (axis cs:0.6,0.6) rectangle (axis cs:0.98,0.98);
	
	

	
		\nextgroupplot
	[
	%xlabel={$x$ (\si{milli\meter})},
	%ylabel={$y$},
	enlargelimits=false,
	%xmin=0.5,
	%xmax=4.5,
	%ymin=0.5,
	%ymax=4.5,
	axis on top,    % <-- (use with `matrix plot*')
	view={0}{90},
	colormap/viridis,
	colorbar horizontal,
	colorbar style={
		at={(0,1.2)},
		anchor=north west,
		xticklabel pos=upper,
		xtick={0,2.5,5,7.5,10},
		xticklabels={0,2.5,5,7.5,10}
		%x tick scale label style={yshift=-25pt, xshift=23pt},
	},
	point meta=explicit,
	point meta min=0.0,
	point meta max=10,
	%xtick={1,2,3,4},
	%ytick={1,2,3,4},
	%xticklabels={1,2,3,4},
	%yticklabels={1,2,3,4},
	]
	\addplot graphics [
	xmin=0, xmax=1,  % Stretch to fit the x-axis limits
	ymin=0, ymax=1    % Stretch to fit the y-axis limits
	] {errorzoom.pdf};
	
	
	\addplot graphics [
	xmin=0.6, xmax=0.98,  % Stretch to fit the x-axis limits
	ymin=0.6, ymax=0.98    % Stretch to fit the y-axis limits
	] {errorsmallbright.pdf};
	
	\draw[fill=none, white] (axis cs:0.6,0.6) rectangle (axis cs:0.98,0.98);
	
	
	
	\end{groupplot}
	\end{tikzpicture}
	
	\caption[Error analysis of simulated diffraction profiles from a pyramid MOT mirror configuration]{Error analysis of simulated diffraction profiles from a pyramid MOT mirror configuration with an aperture of $\SI{0.1}{\milli\meter}\times\SI{0.1}{\milli\meter}$. The methods are based on the Fresnel approximation \textsc{(left)} and direct numerical integration of the Fresnel--Kirchhoff integral \textsc{(middle)}. The error between these methods is shown as a percentage of maximum intensity of the undiffracted beam \textsc{(right)}. The top and bottom rows correspond to simulated areas of $\SI{10}{\milli\meter}\times\SI{10}{\milli\meter}$ and a magnification of the central $\SI{1}{\milli\meter}\times\SI{1}{\milli\meter}$ area, respectively. Insets show additional simulations based on the central area only, highlighting significant boundary artefacts in the direct integration approach.}
	\label{differror}
\end{figure}


\begin{figure}[t]
	\vspace{0em}
	\centering
	\begin{tikzpicture}
	\begin{groupplot}[group style={group size=2 by 2,horizontal sep=0.5cm, vertical sep=0.5cm},height=4cm,x=1.05cm,no markers, y label style={at={(-0.085,0.5)}, anchor=base}]
	
	\nextgroupplot
	[
	ymin=0,
	ymax=0.4,
	%xlabel={$x$ (\si{\milli\meter})},
	ylabel={Error {[\si{\percent}]}},
	xticklabels={,,},
	ytick={0,0.1,0.2,0.3,0.4},
	yticklabels={0,10,20,30,40},
	enlargelimits=false,
	ymajorgrids=true,
	xmajorgrids=true,
	grid style=dashed,
	]
	
	\fill[rainbow7of8, opacity=0.2] (axis cs: 4.5, 0) rectangle (axis cs: 5.5, 2);
	
	\addplot[
	name path=curve,
	color=rainbow3of8,
	] table [x index=0, y index=1, col sep=comma] {E.csv};
	
	
	
	\nextgroupplot
	[
	ymin=0,
	ymax=0.4,
	%xlabel={$x$ (\si{\milli\meter})},
	%ylabel={Error (\si{\percent})},
	xtick={0,1},
	xticklabels={,,},
	ytick={0,0.1,0.2,0.3,0.4},
	yticklabels={,,},
	enlargelimits=false,
	ymajorgrids=true,
	xmajorgrids=true,
	grid style=dashed,
	]
	
	\fill[rainbow7of8, opacity=0.2] (axis cs: 0, 0) rectangle (axis cs: 1, 2);
	
	\addplot[
	name path=curve,
	color=rainbow3of8,
	] table [x index=0, y index=1, col sep=comma] {Es.csv};
	
	\nextgroupplot
	[
	ymin=0,
	ymax=2,
	xlabel={$x$ {[\si{\milli\meter}]}},
	ylabel={Intensity {[a.\,u.]}},
	enlargelimits=false,
	ymajorgrids=true,
	xmajorgrids=true,
	grid style=dashed,
	legend cell align={left},
	legend style={
		at={(1, 1)},  % Adjust (x, y) to move it
		anchor=north east,  % Attach the legend to its top-right corner
		draw=none,           % Optional: Remove border around the legend
		fill=none
	}
	]
	
	\fill[rainbow7of8, opacity=0.2] (axis cs: 4.5, 0) rectangle (axis cs: 5.5, 2);
	
	\addplot[
	name path=curve,
	color=markdiepblauw,
	opacity=1
	] table [x index=0, y index=1, col sep=comma] {Fcs.csv};
	
	\addplot[
	name path=curve,
	color=markrood,
	opacity=0.7
	] table [x index=0, y index=1, col sep=comma] {Kcs.csv};
	
	
	\legend{\footnotesize\hspace{0.3em}Fresnel approx., \footnotesize\hspace{0.3em}direct integration}
	
	\nextgroupplot
	[
	ymin=0,
	ymax=2,
	%xlabel={$x$ (\si{\milli\meter})},
	%ylabel={Norm. intensity},
	enlargelimits=false,
	ymajorgrids=true,
	xmajorgrids=true,
	grid style=dashed,
	yticklabels={,,},
	xtick={0,1},
	]
	
	\fill[rainbow7of8, opacity=0.2] (axis cs: 0, 0) rectangle (axis cs: 1, 2);
	
	\addplot[
	name path=curve,
	color=markdiepblauw,
	opacity=1
	] table [x index=0, y index=1, col sep=comma] {Fscs.csv};
	
	\addplot[
	name path=curve,
	color=markrood,
	opacity=0.7
	] table [x index=0, y index=1, col sep=comma] {Kscs.csv};
	
	\end{groupplot}
	\end{tikzpicture}
	
	\caption[Cross-sectional error analysis of simulated diffraction profiles from a pyramid MOT mirror configuration]{Cross-sectional error analysis of the diffraction profiles from \autoref{differror} based on the Fresnel approximation and direct numerical integration of the Fresnel--Kirchhoff integral. The error \textsc{(top)} between these methods \textsc{(bottom)} is shown as a percentage of maximum intensity of the undiffracted beam. The left and right columns correspond to simulated areas of $\SI{10}{\milli\meter}\times\SI{10}{\milli\meter}$ and $\SI{1}{\milli\meter}\times\SI{1}{\milli\meter}$ (centred), respectively.}
	\label{differrorcross}
\end{figure}


\subsection{Mirror fabrication}


\begin{figure}[t]
	\centering
	\begin{tikzpicture}[remember picture]
		\node[anchor=south west,inner sep=0] (image1) at (0,0) {\includegraphics[width=0.97\textwidth]{{"polish"}.pdf}};
		
		% left four
		\begin{scope}[x={(image1.south east)},y={(image1.north west)}]
			% draw a grid
			%\draw[help lines,xstep=.1,ystep=.1,overlay] (0,0) grid (1,1);
			% draw ticks
			%\foreach \x in {0,1,...,9} { \node [anchor=north,overlay] at (\x/10,0) {0.\x}; }
			%\foreach \y in {0,1,...,9} { \node [anchor=east,overlay] at (0,\y/10) {0.\y}; }
			
			% labels
			%\node[markgrijs, align=center, rotate=0] at (0.0,0.9) {\footnotesize\sc (a)};
			%\node[markgrijs, align=center, rotate=0] at (0.35,0.9) {\footnotesize\sc (b)};
			%\node[markgrijs, align=center, rotate=0] at (0.7,0.9) {\footnotesize\sc (c)};
			
			\node[align=center, rotate=0] at (0.155,-0.05) {\footnotesize\sc (a)};
			\node[align=center, rotate=0] at (0.5,-0.05) {\footnotesize\sc (b)};
			\node[align=center, rotate=0] at (0.845,-0.05) {\footnotesize\sc (c)};
			
		\end{scope}
		
	\end{tikzpicture}
	\caption[Polishing procedure for right-angle prism mirrors]{Polishing procedure for right-angle prism mirrors. \textsc{(a)} The right-angle prism mirror is secured in a custom aluminium mount using a grub screw. A polishing film is attached to a rotating disk of a fibre optic polishing tool. After removal of the required amount of glass from one corner \textsc{(b)}, the prism mirror is flipped inside the mount to polish the adjacent corner \textsc{(c)}. White arrows indicate the direction of polishing.}
	\label{polish} 
\end{figure}



\begin{figure}
	\centering
	\begin{tikzpicture}
		\begin{groupplot}[group style={group size=2 by 1,horizontal sep=2.5cm, vertical sep=2.5cm},height=5.5cm,width=5.5cm,no markers]
			
			
			\nextgroupplot
			[
			x=150,
			y=150,
			ylabel={},
			xmin=0, xmax=1,
			ymin=0, ymax=1,
			xticklabels={,\textcolor{white}{0},},
			yticklabels={,,},
			%legend pos=north west,
			ymajorgrids=true,
			xmajorgrids=true,
			grid style=dashed,
			axis lines=none,        % No axis lines
			%tick style=none,        % No ticks
			%xtick=\empty,           % No x-axis ticks
			%ytick=\empty,           % No y-axis ticks
			xlabel={},              % No x-axis label
			ylabel={},              % No y-axis label
			%enlargelimits=true,     % Keeps the plot within the bounds
			clip=false
			]
			
			\addplot graphics [
			xmin=0, xmax=1,  % Stretch to fit the x-axis limits
			ymin=-0.1, ymax=0.9    % Stretch to fit the y-axis limits
			] {mirrorcut.pdf};
			
			\node[markgrijs, align=center, rotate=0] at (0.23,-0.025) {\footnotesize $5$};
			\node[markgrijs, align=center, rotate=0] at (0.413,0.05) {\footnotesize $b$};
			\node[markgrijs, align=center, rotate=0] at (0.255,0.115) {\footnotesize $a$};
			
			\node[markgrijs, align=center, rotate=0] at (0.93,0.52) {\footnotesize $5$};
			\node[markgrijs, align=center, rotate=0] at (0.832,0.175) {\footnotesize $5$};
			
			\node[markrood, align=center, rotate=0] at (0.57,0.23) {\footnotesize $V$};
			
			
			
			
			\nextgroupplot
			[
			ymin=0,
			ymax=5,
			xmin=0,
			xmax=200,
			ytick={0,1,2,3,4,5},
			xlabel={Polishing time $t$ {[\si{\second}]}},
			ylabel={End facet length $a$ {[\si{\milli\meter}]}},
			ymajorgrids=true,
			xmajorgrids=true,
			grid style=dashed,
			name=plot1,
			]
			\addplot[only marks, markgrijs, error bars/.cd,
			y dir=both, y explicit,
			x dir=both, x explicit] table[x=x, y=y, x error=xe, y error=ye] {
				x     y      xe     ye
				168   0.310  0.1    0.005
				168   0.274  0.1    0.005
				168   0.337  0.1    0.005
				168   0.327  0.1    0.005
				165   0.344  0.1    0.005
				169   0.267  0.1    0.005
				180   0.208  0.1    0.005
			};
			
			\addplot[markrood, thick, domain=0:200, samples=200] {(5 - 2 * ( (6 * 0.01280755 * x)^(1/3) ))};
			
			
			
			%\fill[fill=markdiepblauw, opacity=0.5] (axis cs:163,0.190) rectangle (axis cs:182,0.360);
			
			
			
		\end{groupplot}
		\begin{axis}[
			xticklabel style={rotate=90},
			xlabel style={yshift=0.2cm},
			ylabel style={yshift=-0.1cm},
			xmin=163,
			xmax=182,
			ymin=190,
			ymax=360,
			at={(plot1.north east)},
			xshift=-5pt,
			yshift=-5pt,
			width=3.5cm,
			height=3.5cm,
			anchor=north east,
			xlabel={$t$ {[\si{\second}]}},
			ylabel={$a$ {[\si{\micro\meter}]}},
			]
			\fill[fill=white] (axis cs:163,190) rectangle (axis cs:182,360);
			
		\end{axis}
		\begin{axis}[
			xticklabel style={rotate=90},
			xlabel style={yshift=0.2cm},
			ylabel style={yshift=-0.1cm},
			xmin=163,
			xmax=182,
			ymin=190,
			ymax=360,
			at={(plot1.north east)},
			xshift=-5pt,
			yshift=-5pt,
			width=3.5cm,
			height=3.5cm,
			anchor=north east,
			xlabel={$t$ {[\si{\second}]}},
			ylabel={$a$ {[\si{\micro\meter}]}},
			]
			
			\addplot[only marks, markgrijs, error bars/.cd,
			y dir=both, y explicit,
			x dir=both, x explicit] table[x=x, y=y, x error=xe, y error=ye] {
				x     y      xe     ye
				168   310  0.1    5
				168   274  0.1    5
				168   337  0.1    5
				168   327  0.1    5
				165   344  0.1    5
				169   267  0.1    5
				180   208  0.1    5
			};
			
			\addplot[markrood, thick, domain=160:185, samples=200] {1000*(5 - 2 * ( (6 * 0.01280755 * x)^(1/3) ))};
			
			\draw[markgrijs]       (axis cs:165, 334) -- (axis cs:165, 304) node[below, align=center, markgrijs] {\footnotesize 1};
			\draw[markgrijs]       (axis cs:168+1, 337) -- (axis cs:168+5, 340) node[right, align=center, markgrijs] {\footnotesize 2};
			\draw[markgrijs]       (axis cs:168+1, 327) -- (axis cs:168+7, 315) node[right, align=center, markgrijs] {\footnotesize 3};
			\draw[markgrijs]       (axis cs:168+1, 310) -- (axis cs:168+9, 290) node[right, align=center, markgrijs] {\footnotesize 4};
			
		\end{axis}
	\end{tikzpicture}
	\caption[Right-angle prism mirror geometry and dependence of mirror end facet length on polishing time]{Right-angle prism mirror geometry \textsc{(left)} and dependence of mirror end facet length on polishing time \textsc{(right)}. The schematic shows the substrate dimensions in millimetres, as well as the end facet length $a$ and adjacent length $b$. The red volume $V$ represents the amount of removed material, which can be expressed in terms of the polishing depth $z$ or facet length $a$. On the right, each data point corresponds to a polished substrate (mirrors suitable for assembly are numbered 1--4). The solid line represents a fit based on Preston’s law, which assumes a constant material removal rate ($\dot{V}=\SI{0.0128}{\milli\meter\cubed\per\second}$).}
	\label{polishtime} 
\end{figure}

\begin{figure}[t]
	\centering
	\begin{tikzpicture}[remember picture]
		\node[anchor=south west,inner sep=0] (image1) at (0,0) {\includegraphics[width=0.97\textwidth]{{"mirror"}.pdf}};
		
		% left four
		\begin{scope}[x={(image1.south east)},y={(image1.north west)}]
			% draw a grid
			%\draw[help lines,xstep=.1,ystep=.1,overlay] (0,0) grid (1,1);
			% draw ticks
			%\foreach \x in {0,1,...,9} { \node [anchor=north,overlay] at (\x/10,0) {0.\x}; }
			%\foreach \y in {0,1,...,9} { \node [anchor=east,overlay] at (0,\y/10) {0.\y}; }
			
			% labels
			%\node[markgrijs, align=center, rotate=0] at (0.0,0.9) {\footnotesize\sc (a)};
			%\node[markgrijs, align=center, rotate=0] at (0.35,0.9) {\footnotesize\sc (b)};
			%\node[markgrijs, align=center, rotate=0] at (0.7,0.9) {\footnotesize\sc (c)};
			\node[align=center, rotate=0] at (0.10,0.555) {\footnotesize\sc (a)};
			\node[align=center, rotate=0] at (0.10+0.206,0.555) {\footnotesize\sc (b)};
			\node[align=center, rotate=0] at (0.10+0.412,0.555) {\footnotesize\sc (c)};
			
			\node[align=center, rotate=0] at (0.10,0.055) {\footnotesize\sc (d)};
			\node[align=center, rotate=0] at (0.10+0.206,0.055) {\footnotesize\sc (e)};
			\node[align=center, rotate=0] at (0.10+0.412,0.055) {\footnotesize\sc (f)};
			
			\node[align=center, rotate=0] at (0.815,0.055) {\footnotesize\sc (g)};
			
			
			\node[black, align=center, rotate=0] at (0.514,0.15) {\footnotesize\sc \SI{310}{\micro\meter}};
			\node[black, align=center, rotate=0] at (0.514,0.45) {\footnotesize\sc \SI{344}{\micro\meter}};
			\node[black, align=center, rotate=90] at (0.59,0.29) {\footnotesize\sc \SI{337}{\micro\meter}};
			\node[black, align=center, rotate=90] at (0.44,0.29) {\footnotesize\sc \SI{327}{\micro\meter}};
			
			\node[black, align=center, rotate=0] at (0.514-0.207-0.005,0.15+0.05) {\footnotesize\sc 4};
			\node[black, align=center, rotate=0] at (0.514-0.207-0.005,0.45-0.05) {\footnotesize\sc 1};
			\node[black, align=center, rotate=0] at (0.59-0.207-0.03,0.295) {\footnotesize\sc 2};
			\node[black, align=center, rotate=0] at (0.44-0.207+0.02,0.295) {\footnotesize\sc 3};
			
		\end{scope}
		
	\end{tikzpicture}
	\caption[Fabrication and assembly process of the pyramid MOT mirror system]{Fabrication and assembly process of the pyramid MOT mirror system. \textsc{(a)} Brass collet holding the aluminium mirror mount during machining. \textsc{(b)} Completed mirror holder, which is machined from a gold-coated, aluminium substrate mirror. The gold coating is protected with First Contact during fabrication of the mount. \textsc{(c)} A right-angle prism mirror after its adjacent corners at the base have been polished off. \textsc{(d)} Fully assembled pyramid MOT mirror system. \textsc{(e)} Microscope image showing a top view of the mirrors. The label numbers correspond to those in \autoref{polishtime}. \textsc{(f)} Magnified microscope view of the central aperture, with facet lengths indicated in micrometres. \textsc{(g)} Alignment tool used to position the mirrors within the mount during the glueing process. The alignment is performed via top-view microscope imaging.}
	\label{mirror}
\end{figure}

\subsection{Cold atom detection}

\begin{figure}[t]
	\centering
	\begin{tikzpicture}[remember picture]
		\node[anchor=south west,inner sep=0] (image1) at (0,0) {\includegraphics[width=1\textwidth]{{"pyramidtest"}.pdf}};
		
		% left four
		\begin{scope}[x={(image1.south east)},y={(image1.north west)}]
			% draw a grid
			%\draw[help lines,xstep=.1,ystep=.1,overlay] (0,0) grid (1,1);
			% draw ticks
			%\foreach \x in {0,1,...,9} { \node [anchor=north,overlay] at (\x/10,0) {0.\x}; }
			%\foreach \y in {0,1,...,9} { \node [anchor=east,overlay] at (0,\y/10) {0.\y}; }
			
			% labels
			\node[align=center, rotate=0] at (0.05,0.95) {\footnotesize\sc (a)};
			\node[align=center, rotate=0] at (0.57,0.95) {\footnotesize\sc (b)};
			\node[align=center, rotate=0] at (0.05,0.42) {\footnotesize\sc (c)};
			\node[align=center, rotate=0] at (0.57,0.42) {\footnotesize\sc (d)};
			
		\end{scope}
		
	\end{tikzpicture}
	\caption[Test bed setup for evaluating the standalone performance of a pyramidal MOT before integration with a dipole trap and Fabry--Pérot resonator]{Test bed setup for evaluating the standalone performance of a pyramidal MOT before integration with a dipole trap and Fabry--Pérot resonator. \textsc{(a)} A PTFE holder secures an aluminium mount containing four mirrors arranged in a pyramidal configuration. A cooling beam illuminates the entire pyramid, while a counter-propagating beam plugs the hole at the apex to balance the cooling forces and enable MOT formation. \textsc{(b)} Rear view of the mount, showing the gold-coated surface, which helps reflect additional cooling beams to create an optical molasses at the dipole trap position. \textsc{(c)} The stainless steel mounting structure supports the PTFE holder and aluminium mount, keeping the assembly in place while allowing optical access for the cooling beams. \textsc{(d)} Full assembly inside the glass cell, with anti-Helmholtz coils affixed to a 3D-printed frame, generating the quadrupole magnetic field required for MOT operation.}
	\label{pyramidtest} 
\end{figure}


\begin{figure}[t]
	\vspace{0em}
	\centering
	\begin{subfigure}[b]{.32\linewidth}
		\begin{tikzpicture}
		\centering
		\node[anchor=south west,inner sep=0] (image1) at (0,0) {\includegraphics[width=1.4in] {motbzero.pdf}};
		\begin{scope}[x={(image1.south east)},y={(image1.north west)}]
			\node[white, fill=none, anchor=south west, align=left] at (0,0) {\footnotesize$B=0$};
		\end{scope}
		%\node[color=white] at (1.03,1.1) {\footnotesize \SI{20}{\micro\meter}};
		\end{tikzpicture}
		\centering
		\caption{}
		\label{motbfielda}
	\end{subfigure}
	\begin{subfigure}[b]{.32\linewidth}
		\begin{tikzpicture}
		\centering
		\node[anchor=south west,inner sep=0] (image1) at (0,0) {\includegraphics[width=1.4in] {motbnonzero.pdf}};
		\begin{scope}[x={(image1.south east)},y={(image1.north west)}]
			\node[white, fill=none, anchor=south west, align=left] at (0,0) {\footnotesize$B=\SI{11.3}{\gauss\per\centi\meter}$};
		\end{scope}
		\end{tikzpicture}
		\centering
		\caption{}
		\label{motbfieldb}
	\end{subfigure}
	\begin{subfigure}[b]{.32\linewidth}
		\begin{tikzpicture}
		\centering
		\node[anchor=south west,inner sep=0] (image1) at (0,0) {\includegraphics[width=1.4in] {motbdiff.pdf}};
		\begin{scope}[x={(image1.south east)},y={(image1.north west)}]
			\node[white, fill=none, anchor=south west, align=left] at (0,0) {\footnotesize\sc diff.};
		\end{scope}
		\end{tikzpicture}
		\centering
		\caption{}
		\label{motbfieldc}
	\end{subfigure}
	\caption[Images of a pyramidal mirror configuration with and without an applied magnetic field gradient]{CCD images of a pyramidal mirror configuration installed under ultra-high vacuum (UHV) and continuously illuminated with red-detuned cooling light. \textsc{(a)} Image taken at zero magnetic field ($B = 0$). \textsc{(b)} Image taken at a magnetic field gradient of $B =\SI{11.3}{\gauss\per\centi\meter}$. \textsc{(c)} Difference image obtained by subtracting \textsc{(a)} from \textsc{(b)}.}
	\label{motbfield}
\end{figure}

\section{Conclusion}

\end{document}
